% !TEX root = ../../tree_notes.tex
% =====================================================================
%  Chapter 6 (map v2) -- Constructing Trees from Data      Status: WRITTEN
%  (folder keeps its ch25_constructing_trees name; printed as Chapter 6,
%   label c25.)
%
%  Files, in reading order:
%    00_opener.tex             glance + badges, reading paths, ladder
%                              staircase, the running example T seen through
%                              four data shadows, roadmap, recognition guide,
%                              notation
%    01_parent_edge.tex        6.1 Reconstruct from a Parent Array/Edge List  sec:c25-parent
%    02_degrees_sequences.tex  6.2 Build from Degrees, Sequences and Codes    sec:c25-degree
%    03_is_it_a_tree.tex       6.3 Is It a Tree? Identification & Union-Find  sec:c25-identify
%    04_queries_interactive.tex 6.4 Recover the Tree by Queries (interactive) sec:c25-query
%    99_chapter_end.tex        the 15-part revision contract + Summit (LC 1719)
%
%  This chapter OWNS the reverse move: given raw DATA about a tree (an edge
%  list, a parent array, a degree sequence, an encoding, even answers to
%  queries) REBUILD the tree and PROVE the reconstruction is valid -- exactly
%  one root, n-1 edges, connected, acyclic, sum of degrees 2(n-1). It is the
%  "identify the structure" skill the rest of the book assumes.
%
%  Anti-repetition -- deliberately NOT re-taught here:
%    * the mechanical adjacency build and the parent-guarded rooting DFS
%        -> Chapter 4 (\secref{c04-rooting}, \coderef{c04-root}); we USE them,
%           we add the VALIDITY/root-detection/repair decisions on top
%    * rebuilding from a PAIR of traversals (in+pre, in+post) and serialising
%        -> Chapter 3 (\chref{c03}); this chapter is construction from
%           ARBITRARY data, not from traversal orders
%    * the tin/tout clock, bracket<->tree bijection, distance & LCA machinery
%        -> Chapter 5 (\chref{c24}) and Chapter 11 (\chref{c11}); 6.4 only
%           USES a distance query as an oracle, it does not re-derive it
%    * the full Disjoint-Set-Union structure, path compression
%        -> Chapter 17 (\chref{c28}); we use a union-find PRE-LOOK as the
%           "is it a tree?" test and hand the data structure off there
%    * the full Pruefer-code bijection and Cayley n^(n-2) proof
%        -> Chapter 32 (\chref{c32}); here it is a TEASER that explains why a
%           degree sequence pins a counting problem
%    * trees produced as MINIMAL spanning trees of weighted graphs
%        -> Chapters 17 / 31 (\chref{c28}, \chref{c17})
%  Problems first posed elsewhere are REFERENCED BY NAME, never re-linked
%  (tools/urlsweep.py enforces one URL per problem book-wide).
%
%  Preview one topic:  python tools/preview.py chapters/ch25_constructing_trees/01_parent_edge.tex
% =====================================================================
\bookchapter{Constructing Trees from Data}{c25}
\chapterinfo{LC Easy--Hard}{Written}

% !TEX root = ../../tree_notes.tex
% =====================================================================
%  Chapter 6 opener -- "Chapter at a Glance", reading paths + ladder
%  staircase, the running example T seen through FOUR data shadows
%  (edge list, parent array, degree sequence, queries), roadmap,
%  recognition guide, constraint table, notation.  NO numbered boxes here.
% =====================================================================

\begin{chaptertheorybox}{Chapter at a Glance}
Chapters~4 and~5 pushed \emph{forward}: here is a tree, store it, then lay it on a line.
This chapter pulls \emph{backward}. A contest hands you \term{raw data} -- a list of edges, a
\term{parent array}, a \term{degree sequence}, an \term{encoding}, even \emph{answers to
questions you ask} -- and the real task is to \term{rebuild the tree} and then \emph{prove the
rebuild is legitimate}: connected, acyclic, exactly $n-1$ edges, a single root, degrees summing
to $2(n-1)$. That reverse move -- \emph{identify the structure hiding inside the data} -- is the
skill every later chapter silently assumes the moment it reads \code{edges[i]=[u,v]}. The three
decisions this chapter forces on you:
\begin{itemize}[nosep]
  \item \textbf{Given pairs, is there a tree, and which one?} An edge list or parent array
        determines a tree \emph{only} if it passes the \term{tree certificate}: $n$ nodes,
        $n-1$ edges, connected \emph{and} acyclic (any two of the three force the third). Root
        it by finding the node with no parent. \hfill \lv{2}--\lv{4}
  \item \textbf{Given degrees or a code, can a tree realise them?} A list of positive integers
        is some tree's degree sequence \emph{iff} it sums to $2(n-1)$; the \term{Pr\"ufer code}
        turns every labelled tree into a length-$(n-2)$ sequence, so a degree array even tells
        you \emph{how many} trees exist. \hfill \lv{3}--\lv{5}
  \item \textbf{Given only yes/no and numbers from an oracle, can you still find it?} Ask a
        \emph{distance} or \emph{ancestor} query, get one integer or bit back, and adapt. Each
        honest answer halves the candidates -- reconstructing a hidden tree is binary search on
        a shape. \hfill \lv{2}--\lv{4}
\end{itemize}

\textbf{You will learn to}
\begin{itemize}[nosep]
  \item rebuild an adjacency structure from an edge list or a parent array and locate the root
        in one pass; \hfill\lv{2}--\lv{3}
  \item certify that a graph \emph{is} a tree using the $n-1$-edges / connected / acyclic
        triangle, with a union-find pre-look; \hfill\lv{3}--\lv{4}
  \item decide whether a degree sequence is realisable and construct a tree that realises it;
        \hfill\lv{3}--\lv{4}
  \item read the Pr\"ufer code as the bridge between trees, degree arrays and counting
        (full bijection in \chref{c32}); \hfill\lv{4}--\lv{5}
  \item repair an over-full edge set to a tree (delete the redundant edge) and recover a hidden
        structure from adaptive queries. \hfill\lv{2}--\lv{5}
\end{itemize}

\textbf{Prerequisites.} Representation and rooting from Chapter~4 (\secref{c04-rooting},
\coderef{c04-root}) -- you already know how to pour edges into an \code{adj} list and DFS with a
parent guard; this chapter adds the \emph{validity} and \emph{reconstruction} decisions on top.
The timestamp/distance machinery of Chapter~5 (\chref{c24}) is \emph{used} as an oracle in
\secref{c25-query}, never re-derived.
\end{chaptertheorybox}

% ---------------- Station 1b: reading paths + ladder position ----------------
\begin{chaptertheorybox}{How to Read This Chapter (two paths up the ladder)}
Level badges $\lv1$--$\lv6$ as before: $\lv2$ is ``read the data, rebuild the array'', $\lv3$ is
the solid interview medium, $\lv5$ is an LC Hard reconstruction with a case split.

\textbf{Beginner path ($\lv2$--$\lv3$).} Read \secref{c25-parent} fully -- it turns an edge list
or parent array into a rooted tree and states the tree certificate. Then \secref{c25-identify}
up through the union-find ``is it a tree?'' test. Take every $\lv2$--$\lv3$ solved problem. You
may \emph{skip on a first pass}: the degree/Pr\"ufer material of \secref{c25-degree}, the
interactive protocol of \secref{c25-query}, every $\lv4{+}$ variation, the hardening notes and
the Summit. At the end of this path you can rebuild any tree from an edge list and certify it in
$O(n)$.

\textbf{Full path.} Everything, in order -- it ends with the scaled Summit (LC Hard 1719
\emph{Number of Ways to Reconstruct a Tree}: read an ancestor-pair table, sort nodes by degree,
greedily attach each node to the smallest feasible parent, and decide \emph{unique / many /
impossible}) and the reduction gallery showing that ``rebuild from a parent array'', ``realise a
degree sequence'', ``delete the redundant edge'' and ``find the node by distance queries'' are the
\emph{same} certificate checked against four different data sources. The heavier theory -- the
Disjoint-Set-Union structure itself and the Pr\"ufer/Cayley counting proofs -- lives in
\chref{c28} and \chref{c32}.
\end{chaptertheorybox}

\begin{figure}[H]
\centering
\begin{tikzpicture}[diagram, x=9mm, y=7mm]
  \foreach \st/\lab in {1/lv1,2/lv2,3/lv3,4/lv4,5/lv5,6/lv6} {
    \fill[\lab!22, draw=\lab, line width=0.8pt] (\st-1,0) rectangle (\st,\st*0.52);
    \node[font=\scriptsize\bfseries, text=\lab] at (\st-0.5,\st*0.52-0.42) {L\st};}
  \node[annot, anchor=south west, font=\scriptsize]
    at (0.05,0.62) {6.1 read data, rebuild adj};
  \node[annot, anchor=south west, font=\scriptsize]
    at (1.05,1.14) {6.1/6.3 tree certificate};
  \node[annot, anchor=south west, font=\scriptsize]
    at (2.05,1.66) {6.3 union-find pre-look};
  \node[annot, anchor=south west, font=\scriptsize]
    at (3.05,2.18) {6.2 degree sequence, Pr\"ufer teaser};
  \node[annot, anchor=south west, font=\scriptsize, text=lv5]
    at (4.05,2.70) {6.2 Pr\"ufer counting, 6.4 adaptive queries};
  \node[annot, anchor=south west, font=\scriptsize, text=lv5]
    at (5.05,3.22) {scaled summit (1719)};
\end{tikzpicture}
\caption*{Chapter~6 on the six-step ladder: rebuilding and certifying carry you from raw pairs to
``this is a tree and here is its root'' ($\lv3$); the degree/Pr\"ufer and interactive topics and the
scaled Summit (1719) reach $\lv5$. Beginner path $=$ steps 2--3.}
\end{figure}

% ---------------- The running example ----------------
\subsection*{The Running Example: One Hidden Tree, Four Data Shadows}

This chapter has \emph{one} running example: the general tree $T$ below, drawn once and then
\emph{never shown again} -- instead every topic receives it as a different pile of \emph{data}.
The point is that four unrelated-looking inputs (an edge list, a parent array, a degree sequence,
a bag of distance answers) all describe \emph{this exact tree}, and reconstructing it is the whole
game. Rooted at node $1$, children visited in ascending order:
\[
  \text{edges } \{1\text{--}2,\,1\text{--}3,\,2\text{--}4,\,2\text{--}5,\,3\text{--}6,\,3\text{--}7,\,7\text{--}8\}.
\]

\begin{figure}[H]
\centering
\begin{tikzpicture}[diagram]
  \node[troot] (n1) at (0,0) {$1$};
  \node[tnode] (n2) at (-3,-1.4) {$2$};
  \node[tnode] (n3) at (3,-1.4) {$3$};
  \node[tleaf] (n4) at (-4,-2.8) {$4$};
  \node[tleaf] (n5) at (-2,-2.8) {$5$};
  \node[tleaf] (n6) at (2,-2.8) {$6$};
  \node[tnode] (n7) at (4,-2.8) {$7$};
  \node[tleaf] (n8) at (4,-4.2) {$8$};
  \draw[tedgehl] (n1) -- (n2); \draw[tedgehl] (n1) -- (n3);
  \draw[tedgehl] (n2) -- (n4); \draw[tedgehl] (n2) -- (n5);
  \draw[tedgehl] (n3) -- (n6); \draw[tedgehl] (n3) -- (n7);
  \draw[tedgehl] (n7) -- (n8);
  \node[annot, align=left, font=\scriptsize] at (-6.6,-1.4)
    {$n=8$, $7=n-1$ edges\\ root $=1$\\ leaves $=\{4,5,6,8\}$};
  \node[annot, align=left, font=\scriptsize] at (6.4,-2.0)
    {diameter $4\!\to\!8$\\ is $5$ edges\\ height $=3$};
\end{tikzpicture}
\caption{Running example $T$ (rooted at $1$). Four different data sources -- the edge list above,
the parent array, the degree sequence and a set of distance answers -- each alone is enough to
rebuild this tree, and no two of them look alike. That is the chapter in one picture.}\label{fig:c25-running}
\end{figure}

\medskip
\noindent Keep these four ``shadows'' of $T$ in mind -- they are re-derived, never merely asserted,
in \secref{c25-parent} and \secref{c25-degree}:
\begin{center}
\renewcommand{\arraystretch}{1.25}
\setlength{\tabcolsep}{5pt}
\begin{tabular}{@{}p{0.24\linewidth} p{0.70\linewidth}@{}}
\toprule
\textbf{Data shadow} & \textbf{Value for $T$ (rooted at $1$, nodes $1..8$)} \\
\midrule
edge list (undirected)   & $[1,2],[1,3],[2,4],[2,5],[3,6],[3,7],[7,8]$ -- $n-1=7$ pairs \\
parent array $\parent[\cdot]$ & $[\,0,1,1,2,2,3,3,7\,]$ (index $i$ is node $i$; $\parent[1]=0$ marks the root) \\
degree sequence $\degop[\cdot]$ & $[\,2,3,3,1,1,1,2,1\,]$, sum $=14=2(n-1)$ \\
Pr\"ufer code (teaser)    & $(2,2,1,3,3,7)$ -- length $n-2=6$; node $v$ appears $\degop(v)-1$ times \\
distance answers (6.4)   & $\dist(4,8)=5$, $\dist(5,6)=4$, $\dist(1,8)=3$, \ldots\\
\bottomrule
\end{tabular}
\end{center}

% ---------------- Roadmap diagram ----------------
\begin{figure}[H]
\centering
\begin{tikzpicture}[diagram, node distance=5mm,
    top/.style={boxnode, fill=primaryblue!10, draw=primaryblue, text width=25mm,
                minimum height=11mm, font=\small},
    res/.style={annot, text width=25mm, anchor=north}]
  \node[top] (t1) {\textbf{6.1} Parent / Edge};
  \node[top, right=of t1] (t2) {\textbf{6.2} Degrees \& Codes};
  \node[top, right=of t2] (t3) {\textbf{6.3} Is It a Tree?};
  \node[top, right=of t3] (t4) {\textbf{6.4} By Queries};
  \foreach \a/\b in {t1/t2, t2/t3, t3/t4} \draw[flow] (\a) -- (\b);
  \node[res] at ([yshift=-2mm]t1.south) {rebuild adj; find root; certify $n{-}1$ edges, connected, acyclic};
  \node[res] at ([yshift=-2mm]t2.south) {realisable degree sequence?; greedy build; Pr\"ufer teaser};
  \node[res] at ([yshift=-2mm]t3.south) {union-find pre-look; tree, forest or cycle; drop one edge};
  \node[res] at ([yshift=-2mm]t4.south) {query model; one answer halves it; adaptive search};
\end{tikzpicture}
\caption*{The four topics, left to right, are the \emph{same certificate} tested against four data
sources. 6.1 rebuilds from \emph{pairs}; 6.2 from \emph{degrees/codes}; 6.3 \emph{certifies} (and
repairs) a suspicious edge set; 6.4 recovers the shape from \emph{answers}. Each is construction
from data seen one way.}
\end{figure}

% ---------------- Topic table ----------------
\begin{center}
\renewcommand{\arraystretch}{1.25}
\setlength{\tabcolsep}{4pt}
\begin{tabular}{@{}p{0.24\linewidth} p{0.36\linewidth} c p{0.23\linewidth}@{}}
\toprule
\textbf{Topic} & \textbf{Key results} & \textbf{Levels} & \textbf{Solved problems} \\
\midrule
\secref{c25-parent}\ \ Parent Array / Edge List
  & rebuild adjacency; root-finding; the tree certificate; repair an over-edge set & $\lv2$--$\lv4$
  & \probref{c25-rc2} \\
\addlinespace
\secref{c25-degree}\ \ Degrees, Sequences \& Codes
  & degree sequence realisable $\iff$ sum $2(n{-}1)$; greedy build; Pr\"ufer teaser & $\lv3$--$\lv5$
  & \probref{c25-quad} \\
\addlinespace
\secref{c25-identify}\ \ Is It a Tree? (Union-Find)
  & connected + acyclic + $n{-}1$ edges triangle; union-find test; tree/forest/cyclic & $\lv2$--$\lv3$
  & \probref{c25-rc1} \\
\addlinespace
\secref{c25-query}\ \ Recover by Queries
  & interactive model; what one query reveals; adaptive halving for a hidden node & $\lv2$--$\lv4$
  & \probref{c25-guess} \\
\bottomrule
\end{tabular}
\end{center}

% ---------------- Problem recognition guide ----------------
\subsection*{Problem Recognition Guide --- When You See \emph{this}, Think \emph{that}}

Construction problems announce themselves by \emph{withholding the drawing}: you get pairs,
counts, an array or an oracle instead of pointers. Match the phrasing; every row points at the
section that builds the tool.

\begin{center}
\renewcommand{\arraystretch}{1.3}
\setlength{\tabcolsep}{4pt}
\begin{tabular}{@{}p{0.42\linewidth} p{0.35\linewidth} l@{}}
\toprule
\textbf{When the problem says\ldots} & \textbf{\textcolor{deepblue}{\textbf{Think\ldots}}} & \textbf{Section} \\
\midrule
``\code{edges[i]=[u,v]}'', ``\code{parent[i]} gives the parent of $i$'', a forest given as pairs
  & rebuild \code{adj}, find the root, verify the certificate & \secref{c25-parent} \\
\addlinespace
``is it a valid tree?'', ``one extra edge added -- remove it'', ``which edge is redundant''
  & the $n{-}1$/connected/acyclic test; union-find pre-look & \secref{c25-identify} \\
\addlinespace
``given the degrees of every node, build a tree'', ``how many trees match this degree array''
  & degree sequence sums to $2(n-1)$; Pr\"ufer code counts them & \secref{c25-degree} \\
\addlinespace
``an encoding / grid / nested string -- construct the tree''
  & recursive build-by-rule (quad tree); a stack parse & \secref{c25-degree} \\
\addlinespace
``you may ask \emph{queries}; find the hidden node/edge in $k$ questions''
  & each answer is a bit -- adapt, halve, descend & \secref{c25-query} \\
\addlinespace
``reconstruct so that every listed ancestor-pair holds; is it unique?''
  & sort by degree, greedy parent, transitivity check (Summit) & \probref{c25-summit} \\
\bottomrule
\end{tabular}
\end{center}

\begin{center}
\renewcommand{\arraystretch}{1.25}
\setlength{\tabcolsep}{4pt}
\begin{tabular}{@{}p{0.27\linewidth} p{0.45\linewidth} p{0.20\linewidth}@{}}
\toprule
\textbf{Constraint you see} & \textbf{What it forces} & \textbf{Chapter-6 tool} \\
\midrule
$n-1$ edges promised, graph may be corrupt & do not trust it -- run the certificate/union-find & \secref{c25-identify} \\
a tree \emph{reconstructed} then \emph{queried} many times & rebuild once from the array, answer later & \secref{c25-parent} \\
``how many / which structure fits this degree list'' & Pr\"ufer bijection: degrees pin the count & \secref{c25-degree} (\chref{c32}) \\
interactive judge, tight query budget & information-per-query drives the strategy & \secref{c25-query} \\
$N$ up to $2\cdot10^{5}$ edge list & rebuild \emph{iteratively}; $\bigO{n\alpha}$ union-find & \secref{c25-identify} (\chref{c28}) \\
\bottomrule
\end{tabular}
\end{center}

\begin{notebox}[Not this chapter?]
If the tree is \emph{already given} as pointers or an \code{adj} list and you only need to walk or
store it, that is Chapter~4 (\chref{c04}). If you must rebuild a \emph{binary} tree from two of its
\emph{traversal orders} (inorder+preorder, inorder+postorder) or serialise/parse it, that is
Chapter~3 (\chref{c03}) -- construction from \emph{traversals}, a different data source than the
\emph{arbitrary} data here. If the real question is once the tree stands -- subtree ranges,
distances, LCA -- that is Chapter~5 (\chref{c24}) and Chapter~16 (\chref{c11}). If you need the
Disjoint-Set-Union structure itself (path compression, rank), that is Chapter~17
(\chref{c28}); we borrow only its \emph{pre-look} as a tree test. If you want the full Pr\"ufer
bijection and Cayley's $n^{n-2}$ \emph{proof}, that is Chapter~30 (\chref{c32}). This chapter owns
\emph{reading a tree out of whatever data you were handed, and proving you read it right}.
\end{notebox}

% ---------------- Notation ----------------
\subsection*{Notation Used in This Chapter}

\begin{center}
\renewcommand{\arraystretch}{1.2}
\setlength{\tabcolsep}{4pt}
\begin{tabular}{@{}p{0.20\linewidth} p{0.27\linewidth}@{\quad}p{0.17\linewidth} p{0.27\linewidth}@{}}
\toprule
\textbf{Symbol} & \textbf{Meaning} & \textbf{Symbol} & \textbf{Meaning} \\
\midrule
$T$                 & the running general tree (\figref{c25-running}) & $\parent(v)$ & $v$'s parent; $0$/$-1$ marks the root \\
$\degop(v)$         & number of neighbours of $v$ (degree)           & $\mathcal{E}$ & edge set, $|\mathcal{E}|=n-1$ for a tree \\
$\mathcal{P}$       & the Pr\"ufer code, length $n-2$                & $\mathcal{C}$ & number of connected components \\
$\dist(u,v)$        & edges on the $u$--$v$ path (from \chref{c24})  & $\mathcal{R}$ & a query/answer oracle for a hidden tree \\
$\mathrm{find},\mathrm{uni}$ & union-find ops (\chref{c28})     & $\lv{3}$       & difficulty badge (runs $\lv{1}$--$\lv{6}$) \\
\bottomrule
\end{tabular}
\end{center}

As in earlier chapters, $\dist$ and $\depth$ count \textbf{edges}, general trees are stored as an
adjacency list \code{adj} (\code{vector} of \code{vector<int>}) with a parent guard, and every
listing is complete $\cpp$ (standard header included). Node numbering follows the problem: many
tree-input problems are \textbf{$1$-indexed} (LeetCode 684/685 label nodes $1..n$), which is why
the running example $T$ is $1$-indexed here; \code{adj} arrays are $0$- or $1$-sized to match.

% !TEX root = ../../tree_notes.tex
% =====================================================================
%  6.1  Reconstruct from a Parent Array or Edge List    sec:c25-parent
%  Exports: def:c25-rebuild  thm:c25-cert  thm:c25-degsum  prop:c25-root
%           lst:c25-build  ex:c25-parent-run
%           prob:c25-rc2  lst:c25-rc2  (practice: LC 399)
% =====================================================================
\section{Reconstruct from a Parent Array or an Edge List}\label{sec:c25-parent}

\begin{previouslybox}
\begin{enumerate}[nosep, leftmargin=1.6em]
  \item How do you pour an undirected edge list into an \code{adj} list, and why does the DFS
        need a \emph{parent guard}? (\secref{c04-rooting}, \coderef{c04-root})
  \item For a rooted tree, what single number per node is $0$ for the root and $\ge 1$ for every
        other node, and how do you read it from an array? (\secref{c04-rooting})
  \item How many edges does any tree on $n$ nodes have, and why? (\chref{c01}, \chref{c24})
  \item What is $\sum_v \degop(v)$ for the running tree $T$ of \figref{c25-running}?
\end{enumerate}
\end{previouslybox}

\begin{topicintro}
Almost every tree problem \emph{begs the question in code first}: \code{vector<vector<int>> edges}
or \code{vector<int> parent}. Before any algorithm runs, you must turn that data back into the
object -- an adjacency structure with a distinguished root -- and then \emph{prove the data really
describes a tree}. Chapter~4 showed you how to \emph{store} a tree you already trust; here the
trust must be \emph{earned}. Two objects do all the work. The \term{tree certificate} is the trio
``$n-1$ edges, connected, acyclic'': any two of the three \emph{force} the third, so checking two
is enough. The \term{degree sum} $\sum_v\degop(v)=2(n-1)$ is a cheap arithmetic canary that
instantly rejects corrupt input. Read the parent array's zeros, or an edge list's in-degrees, and
the root falls out. Everything in this section is bookkeeping with one theorem behind it -- and one
Hard problem (LC 685) where the bookkeeping \emph{is} the difficulty.
\textbf{Prerequisites.} \secref{c04-rooting} (adjacency build, parent-guarded rooting,
$\parent/\depth/\size$); \chref{c01} ($n-1$ edges). This section adds \emph{verification and
root detection from raw data}; it does not re-teach traversal.

\textbf{After this section you can}
\begin{itemize}[nosep]
  \item rebuild \code{adj} from an edge list \emph{or} a parent array in one $\bigO{n}$ pass; \hfill\lv{2}
  \item find the unique root from in-degrees or the parent array without any DFS;          \hfill\lv{2}--\lv{3}
  \item certify ``is this data a tree?'' by checking any two of the certificate's three legs; \hfill\lv{3}
  \item repair an \emph{over}-full edge set into a tree by deleting the one redundant edge.  \hfill\lv{4}
\end{itemize}
\end{topicintro}

\begin{recognitionbox}[Pattern card -- rebuild the tree from its data]
\reckey{the input is \emph{not} pointers or a drawing but a pile of \emph{pairs}: \code{edges[i]=[u,v]},
a \code{parent[i]} array, an \code{ancestor[i]} list, or a table of node attributes you must link
into a tree before solving anything else.}
\recuse{
  \item the statement hands you \code{edges} / \code{parent} and the real question (rooting,
        subtree sizes, a repair) needs the \code{adj} built first;
  \item ``is the input a valid tree?'' / ``which edge is extra?'' -- certificate + degree sum;
  \item ``rebuild so a stated rule holds'' (every ancestor pair, every degree) -- \secref{c25-degree}
        and the Summit continue this idea;
  \item the data is \emph{nearly} a tree (one edge too many) -- \probref{c25-rc2}.
}
\recnotuse{
  \item you are given \code{TreeNode*} / \code{root} already -- skip straight to \chref{c03};
  \item you must rebuild from two \emph{traversal orders} -- that is \chref{c03}, not arbitrary data;
  \item the tree is trusted and the work is answering subtree ranges -- \chref{c24}.
}
One-line test: \emph{is there a tree \textbf{hiding inside an array of pairs}? Build it, then
certify it -- this section.}
\end{recognitionbox}

\subsection{Intuition}

An edge list is a tree \emph{in disguise}: $n$ nodes, and a pair $[u,v]$ is a single connection.
To see the object, put every pair into both endpoints' neighbour lists -- exactly the \code{adj}
build of \secref{c04-rooting} -- and, if the tree is \emph{rooted}, orient it with one DFS. The
only \emph{new} questions are: \textbf{which node is the root, and is this even a tree?}

Root detection needs no DFS at all. In a rooted tree exactly one node has no parent. So: from a
\emph{parent array} the root is the index $v$ with $\parent[v]=0$ (or $-1$, or $\parent[v]=v$);
from a directed \emph{edge list} $[u,v]$ meaning ``$u$ is $v$'s parent,'' the root is the node that
never appears as a $v$ (in-degree $0$); from an \emph{undirected} list there is no root until you
\emph{choose} one and root it (\secref{c04-rooting}). Validity is a counting game: $n$ nodes demand
$n-1$ edges, and every node except the root must have exactly one parent.

\begin{figure}[H]
\centering
\begin{tikzpicture}[diagram]
  % highlight the root's slot (label 1) in both rows -- no conditionals needed
  \fill[accentorange!22] (0.62,-0.38) rectangle (1.38,0.38);
  \fill[accentorange!22] (0.62,-1.78) rectangle (1.38,-1.02);
  % parent array row (slots 1..8)
  \node[annot, anchor=east] at (0.4,0) {\small parent};
  \foreach \x/\val in {1/0,2/1,3/1,4/2,5/2,6/3,7/3,8/7} {
    \node[cellnode] at (\x,0) {\val};
  }
  % in-degree row (rooted edge list)
  \node[annot, anchor=east] at (0.4,-1.4) {\small in-deg};
  \foreach \x/\val in {1/0,2/1,3/1,4/1,5/1,6/1,7/1,8/1} {
    \node[cellnode] at (\x,-1.4) {\val};
  }
  \node[annot, align=center] at (4.5,-2.7)
    {slot/label $1$ (orange) is the only $0$ $\Rightarrow$ root $=1$;\\
     in-degrees are a permutation of $\{0,1,\dots,1\}$ -- a $2$ breaks \propref{c25-root}};
\end{tikzpicture}
\caption{Root detection is \emph{arithmetic}, not traversal. Read the parent array: the single $0$
marks node~$1$. Or read the directed edge list's in-degrees: the unique $0$ (node~$1$) is the root
and every other node is exactly~$1$ -- if any node shows in-degree~$2$, the data is \emph{not} a rooted
tree (that is exactly LC 685's alarm, \probref{c25-rc2}).}\label{fig:c25-rootdetect}
\end{figure}

\subsection{Definitions and Notation}

\begin{definition}{Reconstruction from pairs}{c25-rebuild}
Given a set $\mathcal{E}$ of $n-1$ undirected pairs on nodes $\{1,\dots,n\}$, the
\term{reconstruction} is the adjacency structure $\mathit{adj}$ with $v\in\mathit{adj}[u]$ and
$u\in\mathit{adj}[v]$ for every $\{u,v\}\in\mathcal{E}$. A \term{parent array} $\parent[\cdot]$
with $\parent[\mathit{root}]=0$ is reconstructed by adding the directed edge
$\parent[v]\!\to\!v$ for every $v\ne\mathit{root}$. The reconstruction is \emph{valid} if the
result is a tree (rooted, if directed).
\end{definition}

\noindent\textbf{In plain words.} Reconstruction just means ``fill the neighbour lists from the
pairs, and read the root from the array.'' The data can be a lie -- two parents for one node, a
cycle, a disconnected piece -- and validity is what we test next.

\begin{theorem}{The tree certificate}{c25-cert}
For a graph $G$ on $n$ nodes, any \textbf{two} of the following three conditions imply the third,
and all three are equivalent to ``$G$ is a tree'':
\begin{enumerate}[nosep, leftmargin=1.6em]
  \item $G$ has exactly $n-1$ edges;
  \item $G$ is connected;
  \item $G$ is acyclic.
\end{enumerate}
\end{theorem}
\begin{proofbox}
A tree (connected + acyclic) has $n-1$ edges by the classical count (each of the $n-1$ merges in a
spanning process adds one edge). If (1) $n-1$ edges and (2) connected hold but a cycle existed,
breaking one cycle edge keeps connectivity and drops the count to $n-2$, contradicting that a
connected graph needs $\ge n-1$ edges -- so (3). If (1) and (3) hold but $G$ is disconnected with
$k\ge2$ components, each acyclic component on $n_i$ nodes has $n_i-1$ edges, totalling
$n-k<n-1$ edges, contradicting (1) -- so (2). If (2) and (3) hold, a DFS tree uses all $n$ nodes and
every non-tree edge would close a cycle, so exactly $n-1$ edges: (1). Hence any two legs force the
third. \end{proofbox}
\begin{insightbox}
You never check all three -- checking two is the whole certificate, and each check is one cheap
$\bigO{n}$ pass: count the edges ($=n-1$?), and run one DFS or union-find to test connectedness
\emph{or} acyclicity. Pick the two legs your representation makes easiest: an edge list gives
(i) for free; a union-find gives (iii) directly (\secref{c25-identify}).
\end{insightbox}

\begin{theorem}{Handshaking / degree sum for trees}{c25-degsum}
For any tree on $n$ nodes, $\displaystyle\sum_{v} \degop(v) = 2(n-1)$. Consequently the number of
\emph{leaves} (degree~$1$) equals $2+\sum_{v:\degop(v)\ge 2}(\degop(v)-2)$, and a positive-integer
list is a candidate tree degree sequence only if it sums to $2(n-1)$.
\end{theorem}
\begin{proofbox}
Each edge touches exactly two nodes, so counting ``$+1$ per endpoint'' over all edges gives $2$ times
the edge count $=2(n-1)$, which is the sum of degrees. Writing $d-2$ for each node's \emph{excess}
over~$2$ and separating the $\sum(1)=n$ baseline: $\sum_v\degop(v)=n+\sum_v(\degop(v)-1)$, so
$\sum_v(\degop(v)-1)=2(n-1)-n=n-2$; splitting into leaves ($\degop=1$, term $0$) and the rest gives
the leaf identity. \end{proofbox}

\begin{proposition}{Reading the root}{c25-root}
In a rooted tree on $n$ nodes given by a directed edge list $[u,v]$ (``$u$ is $v$'s parent'') or a
parent array: the root is \emph{unique} and is the sole node of in-degree~$0$ (equivalently the sole
index with $\parent[v]=0$), and every other node has in-degree~$1$. A data set with zero or two or
more in-degree-$0$ nodes, or any node of in-degree $\ge2$, is \textbf{not} a rooted tree.
\end{proposition}
\begin{proofbox}
The root has no parent (in-degree $0$); every other node has exactly one parent (in-degree $1$), so
in-degrees are a permutation of $\{0,1,\dots,1\}$ summing to $n-1$. If the root were not unique, two
nodes with no parent would have no common ancestor to connect them (contradicts connectivity), and if
some node had two parents the in-degree sum would exceed $n-1$, breaking \thmref{c25-cert}'s edge
count. Hence exactly one root, and the in-degree histogram certifies the orientation.
\end{proofbox}

\subsection{The algorithm: build \texorpdfstring{$\mathit{adj}$}{adj}, find the root, certify}

\begin{lstlisting}[caption={Rebuild adjacency from an edge list or parent array, find the root, certify validity}, label={lst:c25-build}]
#include <bits/stdc++.h>
using namespace std;

// (A) undirected edge list -> adjacency; root chosen by caller.
vector<vector<int>> buildFromEdges(int n, const vector<vector<int>>& edges) {
    vector<vector<int>> adj(n + 1);
    for (auto& e : edges) { adj[e[0]].push_back(e[1]); adj[e[1]].push_back(e[0]); }
    return adj;
}

// (B) parent array (1-indexed, parent[root]=0) -> children adj + root + validity.
vector<vector<int>> buildFromParent(const vector<int>& par, int& root, bool& ok) {
    int n = (int)par.size() - 1;               // par[0] unused
    vector<vector<int>> adj(n + 1);
    vector<int> indeg(n + 1, 0);
    root = 0; int zeros = 0; ok = true;
    for (int v = 1; v <= n; ++v) {
        int p = par[v];
        if (p == 0) { root = v; ++zeros; }     // candidate root(s)
        else if (p >= 1 && p <= n) { adj[p].push_back(v); ++indeg[v]; }
        else ok = false;                        // parent out of range
    }
    if (zeros != 1) ok = false;                 // need exactly one root
    for (int v = 1; v <= n; ++v)
        if (v != root && indeg[v] != 1) ok = false;   // every other node: one parent
    // optional full check: DFS from root must visit all n nodes (connectivity)
    return adj;
}
\end{lstlisting}
\begin{complexitybox}
Time $\bigO{n}$ -- one pass per edge / array slot, plus (optionally) one DFS for connectivity.
Memory $\bigO{n}$ for \code{adj} and the in-degree histogram. This is the cheapest possible
reconstruction: you never rebuild what a single scan already gives. On a corrupt input the
\code{ok} flag or an in-degree $\ge 2$ is the early warning; \secref{c25-identify} turns that
warning into a full cycle test.
\end{complexitybox}

\subsection{Worked Example}

\begin{example}{Rebuilding $T$ from its two data shadows}{c25-parent-run}
Take the running tree $T$ (\figref{c25-running}), $n=8$.
\emph{From the parent array} $\parent=[0,1,1,2,2,3,3,7]$ (index~$0$ unused): the only $0$ is at
slot~$1$, so $\mathit{root}=1$ and $\mathit{zeros}=1$ (\checkmark). Every other slot has one
parent, so each $v\ne1$ is appended to $\mathit{adj}[\parent[v]]$ and gets in-degree $1$
(\checkmark). \emph{From the edge list} $[1,2],[1,3],[2,4],[2,5],[3,6],[3,7],[7,8]$: seven edges
$=n-1$ (\checkmark), each added to both endpoints. Degree check:
$\degop=[2,3,3,1,1,1,2,1]$, sum $=14=2(n-1)$ (\checkmark). One DFS from $1$ visits all $8$ nodes
(connected). Two legs of \thmref{c25-cert} hold $\Rightarrow$ it \emph{is} a tree, and $1$ is its
root -- reconstructed, not assumed.
\end{example}
\begin{tracebox}[The degree-sum canary rejects two corrupt inputs instantly]
\emph{Corrupt A:} add edge $[4,5]$ to $T$'s edge list. Now $8$ edges for $8$ nodes: count leg (1)
fails ($8\ne n-1$), $\sum_v\degop(v)=16\ne14$ -- reject before any traversal.
\emph{Corrupt B:} set $\parent[7]=2$ (so node $7$ has parents $2$ \emph{and}, via a stray edge,
another). In-degree of $7$ becomes $2$: \propref{c25-root} fires -- no valid root assignment --
reject. This is exactly the shape LC 685 must repair (\probref{c25-rc2}).
\end{tracebox}

\subsection{Variations and Twists}

\begin{itemize}[nosep]
  \item \textbf{Root from a directed edge list}: mark every $v$ that appears as a child; the root
        is the node never marked. $\bigO{n}$ with a boolean array, no histogram needed. \hfill\lv{2}
  \item \textbf{Undirected, root free}: no in-degree tells you the root -- root at \emph{any} node
        (\secref{c04-rooting}) or, if the problem wants the ``true'' root, look for extra clues
        (an \code{isRoot} flag, or node $0$/$1$ by convention). \hfill\lv{2}--\lv{3}
  \item \textbf{Reconstruct-then-query}: build once, store children in \emph{given} order, and
        answer many descendants/subtree questions on the rebuilt tree (\chref{c24} takes over). \hfill\lv{3}
  \item \textbf{Weighted pairs}: the edge $[u,v]$ carries a value; build \code{adj} of
        \code{(to, w)} pairs -- the reconstruction logic is identical (the LC 399 practice below
        is exactly this). \hfill\lv{3}--\lv{4}
  \item \textbf{Near-trees}: $n$ edges (one too many) or a lone in-degree-$2$ node -- do not reject,
        \emph{repair} by deleting the offending edge (\probref{c25-rc2}). \hfill\lv{4}
\end{itemize}

\subsection{Common Pitfalls}\label{sec:c25-parent-pitfalls}

\begin{warningbox}[Trap 1 -- trusting that the input is a tree]
Statements say ``it is a tree'' but corrupt or off-by-one input is a classic WA source. At minimum
check $|\mathcal{E}|=n-1$ and one root (\propref{c25-root}); \secref{c25-identify} adds the full
cycle test. Never run your elegant subtree DP on a graph with a cycle -- the DFS loops forever.
\end{warningbox}
\begin{lstlisting}[caption={\textcolor{warningred}{WRONG}: assume the array is a valid tree, root hard-coded}, language=C++, basicstyle=\ttfamily\small]
int root = 1;                       // what if node 1 isn't the root?
for (int v = 2; v <= n; ++v) adj[par[v]].push_back(v);
\end{lstlisting}
\begin{lstlisting}[caption={\textcolor{successgreen}{RIGHT}: detect the root and validate}, language=C++, basicstyle=\ttfamily\small]
int root = 0, zeros = 0;
for (int v = 1; v <= n; ++v) if (par[v] == 0) { root = v; ++zeros; }
if (zeros != 1) { /* not a valid rooted tree: reject / handle */ }
\end{lstlisting}

\begin{warningbox}[Trap 2 -- mixing 1-indexed data with 0-indexed arrays]
LeetCode tree-input problems label nodes $1..n$ (684, 685) while your \code{adj}/\code{vector}
indices start at $0$. Size arrays $n+1$ and keep the intended index, or translate \emph{consistently}
everywhere -- a half-translation drops node $n$ or reads slot $0$ as a real node.
\end{warningbox}
\begin{lstlisting}[caption={\textcolor{warningred}{WRONG}: $n$-sized array on $1..n$ labels -- index $n$ out of range}, language=C++, basicstyle=\ttfamily\small]
vector<int> indeg(n);       // labels run 1..n but slots are 0..n-1
\end{lstlisting}
\begin{lstlisting}[caption={\textcolor{successgreen}{RIGHT}: size $n+1$, ignore slot $0$}, language=C++, basicstyle=\ttfamily\small]
vector<int> indeg(n + 1, 0);        // one slot per real label
\end{lstlisting}

\begin{warningbox}[Trap 3 -- forgetting the parent guard when treating edges as a tree]
If you rebuild \emph{undirected} \code{adj} and then DFS without a parent guard, the first step back
to the parent re-visits the node you came from and, in a graph that \emph{isn't} yet certified, may
loop. Always pass \code{parent} down the DFS (\coderef{c04-root}) \emph{and} run the certificate.
\end{warningbox}

\subsection{Problem Lab}

The section's anchor is LC 685 -- reconstruct/repair a \emph{rooted} tree given $n$ directed edges
(one too many). It forces every idea above at once: the in-degree histogram, the unique-root
proposition, and the union-find cycle test you meet properly in \secref{c25-identify}.

% ---------------------------------------------------------------------------
\solvedproblem{c25-rc2}{Redundant Connection II}{LeetCode 685}
  {https://leetcode.com/problems/redundant-connection-ii/}{LC Hard L4}
  {directed rooted tree, in-degree (parent) analysis, union-find cycle test, case split}

\begin{problemstatement}[Redundant Connection II]
A \textbf{rooted tree} on $n$ nodes (labels $1\dots n$) had exactly \textbf{one} extra directed edge
added, so the input \code{edges} has $n$ edges (each \code{[u,v]} means ``$u$ is $v$'s parent''). The
extra edge makes either (a) some node have \textbf{two parents}, or (b) a \textbf{cycle}, or both.
Return \emph{one} edge to remove so the rest is a valid rooted tree; if several answers exist,
return the one that appears \textbf{last} in \code{edges}. $3\le n\le 1000$.
\end{problemstatement}

\begin{attemptgate}[Attempt it yourself first -- \lv{4}, give yourself 35 minutes]
A fair first attempt: a rooted tree has $n-1$ edges and every non-root node one parent, so \emph{one}
edge is wrong. Enumerate which edge to delete and test validity each time.
\peek{Look at \textbf{in-degrees}. If some node has in-degree~$2$, one of its two incoming edges is
the culprit -- which one? If \emph{no} node has in-degree~$2$, what does $n$ edges guarantee?}
\nudge{In-degree~$2$ gives two candidate edges. \emph{Temporarily drop the later one} and run a
union-find over the rest: if the rest is still cyclic, the culprit is the \emph{earlier} edge;
otherwise the dropped later edge is the answer. No in-degree~$2$? Just find the edge that closes a
cycle.}
\fullhint{Compute in-degrees. Find a node with in-degree~$2$; label its two incoming edges
\code{first} (earlier index) and \code{banned} (later index). Union-find over all edges, \emph{skipping}
\code{banned}: the first edge whose endpoints are already connected is a cycle -- if a two-parent node
existed return \code{first}, else return that cycle edge. If no cycle remains, return \code{banned}.
Full build and proof below.}
\end{attemptgate}

\solstep{Understanding the Problem}
A valid rooted tree has $n-1$ edges with exactly one node of in-degree $0$ (the root) and $n-1$ nodes
of in-degree $1$. One extra edge pushes the count to $n$ edges and breaks the picture in one of three
\emph{overlapping} ways: a node gains a \textbf{second parent} (in-degree~$2$); or the new edge points
\emph{into} an ancestor, creating a \textbf{cycle}; or it does both at once (the second parent is also
on a cycle). The removal must leave a legal rooted tree, and ties go to the \textbf{last} edge in the
input -- a tie-break you honour by scanning in order and only \emph{proposing} later edges.

\solstep{Key Observations}
\begin{observation}
In-degree~$2$ is \emph{local} evidence: if node $c$ has parents $p_1,p_2$ (edges at indices
$i_1<i_2$), the redundant edge is one of those two -- no other edge can create the second parent. So
the search space collapses from $n$ edges to \textbf{two}.
\end{observation}
\begin{observation}
To choose between the two, \emph{assume} the later edge $i_2$ is the intruder and test whether the
remaining $n-1$ edges (including $i_1$) are \emph{acyclic}. If they are, $i_2$ was indeed safe to
drop. If a cycle survives, the cycle must use $i_1$, so $i_1$ is the true intruder.
\end{observation}
\begin{observation}
With \emph{no} in-degree~$2$ node, $n$ edges on $n$ nodes force exactly one cycle
(\thmref{c25-cert}: edges $>$ $n-1$ $\Rightarrow$ cyclic). A union-find that processes edges in input
order reports the edge that \emph{closes} the cycle -- automatically the last such edge.
\end{observation}

\solstep{Think It Through With Me}
\begin{thinkbox}
\deadend{Try deleting every edge and re-checking ``is it a rooted tree'' with a fresh DFS.}
{$O(n)$ candidates $\times$ $O(n)$ validity $=O(n^{2})$ \emph{tests}, each with fiddly
reachability -- correct, but you reinvent the cycle check $n$ times and still need the
``last edge'' tie-break. A single union-find pass subsumes all of it.}
\deadend{Treat a two-parent node and a cycle with one rule.}{They interact: the extra edge can
\emph{both} give a second parent \emph{and} sit on a cycle. Ignoring the second parent (just chasing
any cycle) returns the wrong edge when the fix is to cut the \emph{earlier} of the two parent edges.}
\aha{Two facts, one scan. In-degree~$2$ names the \emph{only} two candidate edges; a union-find that
\emph{skips the later} candidate tells you which of the two survives the cycle test. No two-parent
node means the cycle edge itself is the answer.}
\end{thinkbox}

\solstep{Solution}
Compute in-degrees. If a node $c$ has two incoming edges at input indices $i_1<i_2$, set
\code{banned}$=i_2$ and \code{first}$=i_1$. Run union-find over $1..n$ processing every edge
\emph{except} \code{banned} in input order: on the first edge whose endpoints are already in one set,
a cycle exists -- return edges[\code{first}] if a two-parent node was found (dropping the later edge
did \emph{not} kill the cycle), otherwise return that closing edge. If the scan finishes with no
cycle, return edges[\code{banned}] (the later parent edge was the redundant one) -- or an empty
answer in the degenerate no-cycle-no-two-parent case, which cannot occur for valid input.

\solstep{Proof of Correctness}
\begin{proofbox}
\emph{Case 1: some node has in-degree~$2$.} Let its parent edges be $e_1$ (earlier) and $e_2$ (later).
Any valid tree must drop at least one of them, since no node may have two parents.
If removing $e_2$ yields an acyclic edge set, those $n-1$ edges with the surviving single-parent
structure are a rooted tree, so $e_2$ is a correct removal; $e_2$ is later than $e_1$, matching the
tie-break, so we return $e_2$.
If removing $e_2$ still leaves a cycle, then $e_1$ lies on a cycle (the union-find reports a cycle
among the edges excluding $e_2$), so removing $e_2$ alone cannot work -- the redundant edge must be
$e_1$, which also resolves the double parent. Hence $e_1$ is returned.
\emph{Case 2: no node has in-degree~$2$} and $n$ edges exist: exactly one cycle ($\thmref{c25-cert}$
contrapositive). The union-find returns the last edge closing that cycle. Removing it breaks the only
cycle and keeps every in-degree $\le1$, so the rest is a rooted tree; it is the latest valid answer
because we scan in order and any earlier cycle edge is not chosen. In both cases the returned edge is
exactly the one whose removal restores a valid rooted tree, chosen to appear last among ties.
\end{proofbox}

\solstep{Implementation}
\begin{lstlisting}[caption={Redundant Connection II -- in-degree scan + union-find with one edge banned (LeetCode 685)}, label={lst:c25-rc2}]
#include <bits/stdc++.h>
using namespace std;

class Solution {
    vector<int> p;
    int find(int x) { return p[x] == x ? x : p[x] = find(p[x]); }
    bool uni(int a, int b) {                    // false == already connected (cycle)
        a = find(a); b = find(b);
        if (a == b) return false;
        p[b] = a; return true;
    }
public:
    vector<int> findRedundantDirectedConnection(vector<vector<int>>& edges) {
        int n = (int)edges.size();              // n nodes, n edges
        p.assign(n + 1, 0);
        for (int i = 0; i <= n; ++i) p[i] = i;
        vector<int> indeg(n + 1, 0);
        for (auto& e : edges) ++indeg[e[1]];

        int banned = -1, first = -1;            // the two edges into an in-degree-2 node
        for (int i = 0; i < n; ++i) if (indeg[edges[i][1]] == 2) {
            if (first == -1) first = i;         // earlier incoming edge
            else             banned = i;        // later incoming edge -> skip it first
        }
        for (int i = 0; i < n; ++i) {
            if (i == banned) continue;          // pretend e_later is already removed
            if (!uni(edges[i][0], edges[i][1])) // cycle closes here
                return (banned != -1) ? edges[first] : edges[i];
        }
        return (banned != -1) ? edges[banned] : vector<int>{};
    }
};
\end{lstlisting}

\solstep{Trace}
\begin{tracebox}[Two official inputs, one per case]
\emph{Ex.~1} \code{edges=[[1,2],[1,3],[2,3]]}. In-degrees: node~$3$ has $2$ (from~$1$ and~$2$). So
\code{first=1} ($[1,3]$), \code{banned=2} ($[2,3]$). Union-find skipping index~$2$: union $1,2$;
union $1,3$; no cycle. Loop ends clean $\Rightarrow$ return edges[\code{banned}]$=[2,3]$. \checkmark\\
\emph{Ex.~2} \code{edges=[[1,2],[2,3],[3,4],[4,1],[1,5]]}. All in-degrees $\le1$ (node~$1$ from~$4$),
so \code{banned=first=-1}. Union in order: $1\text{–}2$, $2\text{–}3$, $3\text{–}4$ merge $\{1,2,3,4\}$;
edge $[4,1]$: $4,1$ already connected $\Rightarrow$ cycle, \code{banned=-1} so return $[4,1]$.
\checkmark (matches the expected $[4,1]$).
\end{tracebox}

\solstep{Complexity}
\begin{complexitybox}
Time $\bigO{n\,\alpha(n)}$ -- one in-degree pass and one union-find pass; $\alpha$ is the inverse
Ackermann ($\le 4$ for all realistic $n$). Space $\bigO{n}$. At $n\le1000$ the $\bigO{n^{2}}$ brute
force also passes, but this is the clean linear solution. Verdict: \ac.
\end{complexitybox}

\solstep{Pitfalls and Variants}
\begin{warningbox}
Three ways to break it: (i) skipping the \emph{earlier} parent edge instead of the later one -- the
tie-break and the cycle logic both want the \emph{later} edge banned first. (ii) Running the
union-find \emph{including} the banned edge -- then the two-parent node always looks cyclic and you
return the wrong parent. (iii) Forgetting $1$-indexing: nodes are $1..n$, so \code{p}/\code{indeg}
are sized $n+1$. \emph{Variant:} LC 684 is the undirected, single-cycle special case (just
Case~2) -- \probref{c25-rc1}.
\end{warningbox}

\begin{insightbox}[Key takeaways]
\begin{enumerate}[nosep, leftmargin=1.6em]
  \item \textbf{In-degree is evidence.} A second parent pins the culprit to two edges -- the whole
        search space -- because the degree sum \thmref{c25-degsum} and \propref{c25-root} leave no
        room for a third possibility.
  \item \textbf{One union-find pass replaces $n$ validity tests} and hands you the
        ``appears last'' edge for free, because edges are consumed in input order.
  \item \textbf{Reconstruction and repair are the same skill:} read structure out of degrees and
        pairs, then \emph{test} it -- exactly the certificate of \thmref{c25-cert}.
\end{enumerate}
\end{insightbox}

\begin{hardeningbox}
\begin{itemize}[nosep, leftmargin=1.4em]
  \item \textbf{Both defects at once}: if a two-parent node is \emph{also} on the cycle, banning the
        later edge still leaves a cycle, and the code correctly returns the \emph{earlier} parent edge
        -- worth an explicit test input like \code{[[2,1],[3,1],[4,2],[1,4]]}.
  \item \textbf{Larger $n$}: the path-compressed union-find keeps it near-linear; avoid a per-candidate
        DFS re-check (the $O(n^2)$ trap). Full DSU theory (rank, compression) is \chref{c28}.
  \item \textbf{Generalisation}: ``delete $k$ edges to make a forest'' needs the component count of
        \secref{c25-identify}; this problem is the $k=1$, rooted, tie-broken special case.
\end{itemize}
\end{hardeningbox}

\subsection{Practice Problems}

\begin{stuckbox}
\begin{itemize}[nosep, leftmargin=1.4em]
  \item Did you detect the root by \emph{counting} (one in-degree-$0$ node), not by hard-coding?
        (Trap 1, \propref{c25-root}.)
  \item Are your arrays sized for $1$-indexed labels? (Trap 2.)
  \item When choosing between two parent edges, did you \emph{ban the later} one first and test for a
        surviving cycle? (\probref{c25-rc2}.)
  \item Did you use the $n-1$ / degree-sum canary before trusting the data? (\thmref{c25-degsum}.)
\end{itemize}
\end{stuckbox}

\begin{practiceset}
  \practice{Evaluate Division}{LeetCode 399}
    {https://leetcode.com/problems/evaluate-division/}{LC Medium L3}
    {You are handed raw \emph{pairs} \code{equations[i]=[a,b]} with a value $a/b$, and must answer
    queries $x/y$. There is no tree given -- you must \emph{reconstruct} a structure from the pairs:
    build an adjacency list where $a\!\to\!b$ carries weight $w$ and $b\!\to\!a$ carries $1/w$, then
    each query is a weighted path product through that rebuilt graph (DFS/BFS, or a weighted
    union-find). Recognise it as ``build \code{adj} from pairs, then answer'' -- \secref{c25-parent}
    in weighted clothing, with connectivity deciding ``answerable vs $-1$''.}
\end{practiceset}

\begin{recallbox}[10-second refresher -- 6.1 Reconstruct from Parent Array / Edge List]
Turn pairs into \code{adj} in one pass; find the root as the sole in-degree-$0$ node
(\propref{c25-root}); certify with any \emph{two} of ``$n-1$ edges, connected, acyclic''
(\thmref{c25-cert}); reject bad input with $\sum_v\degop(v)=2(n-1)$ (\thmref{c25-degsum}). To \emph{repair}
one extra edge (LC 685), let the in-degree-$2$ node pin two candidates and a union-find (skipping the
later) choose (\probref{c25-rc2}).
\selftest{
  \item Rebuild $T$ from its parent array $[0,1,1,2,2,3,3,7]$: name the root and the leaves.
  \item State the tree certificate and prove that ``$n-1$ edges $+$ acyclic'' forces connectedness.
  \item Why must a rooted tree's in-degrees be a permutation of $\{0,1,\dots,1\}$?
  \item In LC 685, why do you ban the \emph{later} parent edge, not the earlier, before the cycle test?
  \item Give the degree-sum of any tree on $n$ nodes and explain the leaf formula.
}
\end{recallbox}

% !TEX root = ../../tree_notes.tex
% =====================================================================
%  6.2  Build from Degrees, Sequences and Codes       sec:c25-degree
%  Exports: def:c25-degseq  thm:c25-degreal  def:c25-prufer  prop:c25-prufdeg
%           lst:c25-degbld  ex:c25-deg-run  def:c25-quad  lst:c25-quad
%           prob:c25-quad  lst:c25-quad  (practice: LC 210)
% =====================================================================
\section{Build from Degrees, Sequences and Codes}\label{sec:c25-degree}

\begin{previouslybox}
\begin{enumerate}[nosep, leftmargin=1.6em]
  \item State the handshaking identity for a tree's degrees and the exact value of
        $\sum_v\degop(v)$ for the running tree $T$. (\thmref{c25-degsum})
  \item Why can a graph with a node of in-degree~$2$ never be a rooted tree?
        (\propref{c25-root})
  \item How many edges must a realisation of a degree sequence have, and how do you
        read that straight off the sequence? (\thmref{c25-cert})
  \item In the leaf-peel of a tree, why is there \emph{always} a node of degree~$1$ to remove?
        (\chref{c01})
\end{enumerate}
\end{previouslybox}

\begin{topicintro}
Section~6.1 rebuilt a tree from \emph{pairs}. Here the input is even coarser: a list of
\emph{numbers}. You are handed a \term{degree sequence} $d_1,\dots,d_n$ and asked (i) does
\emph{any} tree have exactly these degrees, and (ii) if so, \emph{build one} -- or even
\emph{count} them. The answer is shockingly clean: a list of positive integers is a tree degree
sequence \textbf{iff} it sums to $2(n-1)$ (\thmref{c25-degreal}), and the constructive proof is a
leaf-peel that is the exact inverse of the \term{Pr\"ufer code} -- the length-$(n-2)$ sequence
that bijects labelled trees with multisets, and the reason a degree array also tells you
\emph{how many} trees realise it (\propref{c25-prufdeg}; the full bijection lives in
\chref{c32}). The same ``rebuild a branching object from aggregate data'' instinct reappears in a
completely different guise: \term{constructing a quad tree} by recursively partitioning a grid
(LC 427), where the data is a matrix and the tree is a divide-and-conquer shape. And a
\emph{topological order} (LC 210) is a linear ``tree'' peeled from in-degrees -- the identical
degree-peel mechanic. One idea -- \emph{read structure back out of degrees} -- three faces.
\textbf{Prerequisites.} \thmref{c25-degsum}, \propref{c25-root} from \secref{c25-parent}. The
Pr\"ufer decoding is \emph{used} as a construction; its bijective proof and Cayley's count are
\chref{c32}. This section adds \emph{realisability, greedy construction and counting from a
degree array}, plus recursive-shape construction.

\textbf{After this section you can}
\begin{itemize}[nosep]
  \item decide in $O(n)$ whether a degree list is realisable by a tree, and say why;         \hfill\lv{3}
  \item construct \emph{an} explicit tree realising a given degree sequence by a leaf-peel;   \hfill\lv{3}--\lv{4}
  \item read the Pr\"ufer code as $\degop(v)-1$ occurrences and count realising trees;        \hfill\lv{4}--\lv{5}
  \item build a quad tree (or any recursive partition tree) from raw grid data by divide and conquer. \hfill\lv{3}
\end{itemize}
\end{topicintro}

\begin{recognitionbox}[Pattern card -- realise a tree from degrees / codes / aggregate data]
\reckey{the input is a list of \emph{numbers}, not pairs: a \code{degree[]} array, an
\code{ancestor-count}, a Pr\"ufer-like \emph{code}, a grid you must turn into a partition tree,
or an in-degree table you peel into an order.}
\recuse{
  \item ``does a tree exist with these degrees? / build one'' -- \thmref{c25-degreal} plus the leaf-peel;
  \item ``how many labelled trees have this degree sequence?'' -- Pr\"ufer multinomial \propref{c25-prufdeg};
  \item ``compress this $n\times n$ grid into a recursive tree'' -- quad tree, \probref{c25-quad};
  \item ``order the tasks so every dependency comes first'' -- in-degree peel (LC 210, the practice).
}
\recnotuse{
  \item you already have the edges and only need the root -- \secref{c25-parent};
  \item the sequence is an \emph{arbitrary graph}'s degrees (not necessarily a tree): that is the
        Erd\H{o}s--Gallai / Havel--Hakimi question for general graphs, out of scope here;
  \item you must answer subtree/path queries on a trusted tree -- \chref{c24}.
}
One-line test: \emph{is the input a pile of \textbf{numbers} that must become a
\textbf{shape}? Check the sum, then peel leaves -- this section.}
\end{recognitionbox}

\subsection{Intuition}

Forget the labels for a moment. Every tree on $n$ nodes has $n-1$ edges, and each edge donates
$1$ to two endpoints, so the degrees must sum to $2(n-1)$ -- there is simply no other total
(\thmref{c25-degsum}). The surprising direction is the \emph{converse}: as long as every $d_i\ge1$
and the sum is $2(n-1)$, you can \emph{always} wire up a tree with exactly those degrees. The
recipe is a peeling game. Because the average degree is $2(n-1)/n<2$, some node has degree~$1$ --
a forced \emph{leaf}. Pin that leaf onto any node that still has degree to spare ($>1$), spend one
unit from each, and recurse on the smaller problem. Each pin removes exactly one leaf and never
creates an illegal degree, so after $n-1$ pins you have a genuine tree. This is the same bookkeeping
that underlies Pr\"ufer encoding, seen from the decode side.

\begin{figure}[H]
\centering
\begin{tikzpicture}[diagram]
  % degree row for T (nodes 1..8)
  \node[annot, anchor=east] at (0,0) {\small $\degop$};
  \foreach \x/\val in {1/2,2/3,3/3,4/1,5/1,6/1,7/2,8/1} {
    \node[cellnode] at (\x,0) {\val};
  }
  % highlight the leaves (deg 1): 4,5,6,8
  \fill[successgreen!20] (3.62,-0.38) rectangle (4.38,0.38);
  \fill[successgreen!20] (4.62,-0.38) rectangle (5.38,0.38);
  \fill[successgreen!20] (5.62,-0.38) rectangle (6.38,0.38);
  \fill[successgreen!20] (7.62,-0.38) rectangle (8.38,0.38);
  % deg-1 row (Prufer multiplicity)
  \node[annot, anchor=east] at (0,-1.4) {\small $\degop{-}1$};
  \foreach \x/\val in {1/1,2/2,3/2,4/0,5/0,6/0,7/1,8/0} {
    \node[cellnode] at (\x,-1.4) {\val};
  }
  \node[annot, align=center] at (4.5,-2.8)
    {top row sums to $14=2(n-1)$ \checkmark\ realisable;\\
     green cells are the forced leaves the peel starts with;\\
     bottom row sums to $6=n-2$ and \emph{is} the Pr\"ufer multiset of $T$};
\end{tikzpicture}
\caption{Two arithmetic views of the same degree array. The \emph{top} row is $T$'s degrees; the
$\sum=2(n-1)$ test is the whole realisability certificate (\thmref{c25-degreal}). The \emph{bottom}
row subtracts $1$ from each entry -- its sum is $n-2$, and reading node $v$ exactly $\degop(v)-1$
times gives the Pr\"ufer \emph{multiset} $(2,2,1,3,3,7)$ of \secref{c25-parent}'s teaser.}\label{fig:c25-degrow}
\end{figure}

\subsection{Definitions and Notation}

\begin{definition}{Tree degree sequence; realisability}{c25-degseq}
A list $d=(d_1,\dots,d_n)$ of positive integers is \term{tree-realisable} if there exists a tree
$G$ on nodes $\{1,\dots,n\}$ with $\degop(i)=d_i$ for every $i$. The construction problem asks for
one such $G$; the counting problem asks how many labelled $G$ realise $d$.
\end{definition}

\noindent\textbf{In plain words.} ``Can these numbers be the degrees of \emph{some} tree, and if so
show me the tree (or count them).'' Nothing about positions or edges is given -- only the per-node
totals.

\begin{theorem}{Degree-sequence realisability for trees}{c25-degreal}
A list $d_1,\dots,d_n$ of positive integers is tree-realisable \textbf{if and only if}
$\displaystyle\sum_{i=1}^{n} d_i = 2(n-1)$. When $n\ge2$, the leaf-peel of \coderef{c25-degbld}
constructs a realising tree.
\end{theorem}
\begin{proofbox}
\emph{(Only if)} Every tree has $n-1$ edges and each edge contributes $2$ to the degree sum
(\thmref{c25-degsum}), so any realisation forces $\sum d_i=2(n-1)$, and clearly each $d_i\ge1$
(a node in a connected graph on $n\ge2$ vertices has a neighbour).
\emph{(If)} We run the peel. Since every $d_i\ge1$ and $\sum d_i=2(n-1)<2n$, the average is below
$2$, so some node has $d_i=1$ -- a leaf exists at every stage. Pin a current leaf~$\ell$ to any node
$p$ of current degree $>1$ and append edge $(\ell,p)$, decrementing both. The invariant ``all
remaining degrees positive and their sum equals twice the edges still to add'' is preserved: the sum
drops by $2$ (one from $\ell$, one from $p$) and we remove one leaf, so there are $n-1$ iterations
before all degrees reach~$0$. The $n-1$ produced edges connect all $n$ nodes (each new edge ties a
not-yet-used leaf into the growing component, and the final edge joins the last two leaves), giving
$n-1$ edges and, by \thmref{c25-cert}, a tree whose degrees are exactly $d$. \end{proofbox}
\begin{insightbox}
This is a \emph{sharp} contrast with general graphs: for an arbitrary graph the realizability test
is the Erd\H{o}s--Gallai inequalities or a Havel--Gakimi peel. For \emph{trees} the single equation
$\sum d_i=2(n-1)$ (plus $d_i\ge1$) is necessary \emph{and} sufficient -- trees are rigid enough that
one number decides everything. The degree array is ``missing information'' only about \emph{which}
tree, never about \emph{whether} one exists.
\end{insightbox}

\begin{definition}{Pr\"ufer code (construction view)}{c25-prufer}
The \term{Pr\"ufer code} of a labelled tree on $\{1,\dots,n\}$ is a sequence of length $n-2$ over
the alphabet $\{1,\dots,n\}$ obtained by repeatedly deleting the smallest leaf and recording its
neighbour. Crucially, $\degop(v)-1$ equals the number of times $v$ appears in the code, so
\emph{any} ordering with those multiplicities decodes to \emph{a} tree realising $d$. (The
deletion--recording bijection and its inverse are proven in full in \chref{c32}; here we only
\emph{use} the decode.)
\end{definition}

\begin{proposition}{Counting realising trees from a degree sequence}{c25-prufdeg}
If $d_1,\dots,d_n$ is tree-realisable ($n\ge2$), the number of \emph{labelled} trees whose degree
sequence is exactly $d$ is
\[
  \frac{(n-2)!}{\prod_{i=1}^{n}(d_i-1)!}.
\]
Summing this multinomial over all realisable $d$ recovers Cayley's $n^{\,n-2}$ labelled trees.
\end{proposition}
\begin{proofbox}
By \defref{c25-prufer}, a labelled tree is uniquely coded by a length-$(n-2)$ sequence in which
symbol $i$ appears exactly $d_i-1$ times; the sum of these multiplicities is
$\sum(d_i-1)=2(n-1)-n=n-2$, so the count is the number of distinct permutations of that multiset --
the multinomial $(n-2)!/\prod(d_i-1)!$. Since the Pr\"ufer map is a bijection (proven in
\chref{c32}), this multinomial \emph{is} the number of realising labelled trees. \end{proofbox}

\begin{example}{How many trees look like $T$'s degrees?}{c25-deg-count}
$d=[2,3,3,1,1,1,2,1]$, so $d_i-1=[1,2,2,0,0,0,1,0]$ and the multinomial is
$6!/(1!\,2!\,2!\,0!\,0!\,0!\,1!\,0!)=720/4=\mathbf{180}$ distinct labelled trees share $T$'s degree
sequence. Cayley gives $8^{6}=262144$ trees in total on $8$ labels; the $180$ above are just those
with this particular degree profile -- the array narrows the world by three orders of magnitude, yet
still cannot name \emph{which} tree. \end{example}

\subsection{The algorithm: peel leaves to realise a degree sequence}

\begin{lstlisting}[caption={Construct a tree realising a degree sequence, or report it impossible}, label={lst:c25-degbld}]
#include <bits/stdc++.h>
using namespace std;

// vertices are 0..n-1; deg[v] = required degree of v.
// Returns n-1 edges realising deg, or an empty list if deg is not tree-realisable.
vector<pair<int,int>> treeFromDegrees(const vector<int>& deg) {
    int n = (int)deg.size();
    long long sum = 0;
    for (int d : deg) { if (d < 1) return {}; sum += d; }     // positive required
    if (n == 1) return {};                                   // single node, no edges
    if (sum != 2LL * (n - 1)) return {};                     // the WHOLE test (Thm 6.2)

    vector<int> cur = deg;                                   // live degrees during the peel
    // Pruefer-style leaf-min-heap: always peel the smallest current leaf.
    priority_queue<int, vector<int>, greater<int>> leaf;
    for (int i = 0; i < n; ++i) if (cur[i] == 1) leaf.push(i);

    // Order vertices by (deg-1) copies -> a Pruefer sequence of length n-2.
    vector<int> prufer;
    for (int i = 0; i < n; ++i) for (int k = 1; k < deg[i]; ++k) prufer.push_back(i);

    vector<pair<int,int>> E;
    for (int x : prufer) {                                   // decode: attach a leaf to x
        int l = leaf.top(); leaf.pop();
        E.emplace_back(l, x);
        if (--cur[x] == 1) leaf.push(x);                     // x became a leaf
    }
    int a = leaf.top(); leaf.pop();                          // last two leaves join
    int b = leaf.top(); leaf.pop();
    E.emplace_back(a, b);
    return E;                                                // exactly n-1 edges, degrees = deg
}
\end{lstlisting}
\begin{complexitybox}
Time $\bigO{n\log n}$ -- the Pr\"ufer multiset holds $n-2$ entries and every heap operation is
$\log n$; building the multiset is itself $O(n)$ since $\sum(\degop(v)-1)=n-2$. Memory $\bigO{n}$.
The realisability \emph{decision} alone is $O(n)$ (one sum). The construction is exactly the
Pr\"ufer \emph{decoder} run on a conveniently chosen code, which is why it always succeeds when
\thmref{c25-degreal}'s equation holds.
\end{complexitybox}

\subsection{Worked Example}

\begin{example}{Realising $T$'s degree sequence by the peel}{c25-deg-run}
Feed $d=[2,3,3,1,1,1,2,1]$ (nodes $0..7$ relabel $1..8$) to \coderef{c25-degbld}. The sum is $14$
$=2(n-1)$ \checkmark, so we proceed. Initial leaves (degree~$1$): $\{4,5,6,8\}$ (1-indexed). The
Pr\"ufer multiset, node $v$ listed $\degop(v)-1$ times, is $\{1,\,2,2,\,3,3,\,7\}$ (length~$6$).
Decoding in ascending order: attach leaf $4\!\to\!1$ (node~$1$ drops to~$1$, joins the leaf pool);
attach leaf $5\!\to\!2$; attach leaf $6\!\to\!2$ (node~$2$ now degree~$1$, a leaf); attach leaf
$1\!\to\!3$; attach leaf $2\!\to\!3$ (node~$3$ drops to~$1$); then node~$7$; the final two free
leaves join. The result is \emph{some} tree with degrees exactly $[2,3,3,1,1,1,2,1]$ -- not
necessarily $T$ itself, but a legal realisation, which is all the theorem promises. Re-running with a
different Pr\"ufer \emph{order} would yield a different one of the $180$ trees counted in
\exref{c25-deg-count}.
\end{example}
\begin{tracebox}[The peel never gets stuck]
At every step the ``positive degrees, sum $=$ twice-remaining-edges'' invariant holds, so the
average degree of the live graph is below $2$ and a degree-$1$ node is guaranteed to be in the heap.
If you \emph{do} hit an empty leaf heap, the input failed \thmref{c25-degreal} -- e.g.\ $d=[2,2,2]$
sums to $6\ne4$, or $d=[3,1]$ has a node wanting $3$ with only one partner: both return empty.
\end{tracebox}

\subsection{From grids to recursive trees: the quad tree}

The degree array is aggregate data of type ``counts''. A \term{quad tree} is aggregate data of type
``a square matrix'', and the reconstruction is \emph{recursive partition}: if a square is uniform,
it collapses to a leaf; otherwise split it into four quadrants and rebuild each. The output tree's
shape is forced by the grid -- no edges are ever given, yet a genuine tree pops out. This is the
construction soul of LC 427.

\begin{figure}[H]
\centering
\begin{tikzpicture}[diagram, scale=0.75, transform shape]
  % ---- 4x4 grid on the left ----
  \begin{scope}[xshift=-3.2cm]
    \node[annot] at (1,-0.4) {\small input grid};
    \foreach \i in {0,1,2,3} \foreach \j in {0,1,2,3} {
      \draw[draw=primaryblue!60, line width=0.6pt] (\i,-\j) rectangle (\i+1,-\j-1);
    }
    % shade the 1-cells: bottom-left 2x2 all 1, plus top-left single
    \fill[accentorange!35] (0,-2) rectangle (2,-4);   % bottom-left quadrant = 1
    \node[font=\scriptsize] at (0.5,-2.5) {$1$}; \node[font=\scriptsize] at (1.5,-2.5) {$1$};
    \node[font=\scriptsize] at (0.5,-3.5) {$1$}; \node[font=\scriptsize] at (1.5,-3.5) {$1$};
    \node[font=\scriptsize] at (0.5,-0.5) {$0$}; \node[font=\scriptsize] at (1.5,-0.5) {$0$};
    \node[font=\scriptsize] at (2.5,-0.5) {$0$}; \node[font=\scriptsize] at (3.5,-0.5) {$0$};
    \node[font=\scriptsize] at (0.5,-1.5) {$0$}; \node[font=\scriptsize] at (1.5,-1.5) {$0$};
    \node[font=\scriptsize] at (2.5,-1.5) {$0$}; \node[font=\scriptsize] at (3.5,-1.5) {$0$};
    \node[font=\scriptsize] at (2.5,-2.5) {$1$}; \node[font=\scriptsize] at (3.5,-2.5) {$0$};
    \node[font=\scriptsize] at (2.5,-3.5) {$0$}; \node[font=\scriptsize] at (3.5,-3.5) {$0$};
    \draw[draw=warningred, line width=1pt, dashed] (2,-0) -- (2,-4);
    \draw[draw=warningred, line width=1pt, dashed] (0,-2) -- (4,-2);
  \end{scope}
  % ---- quad tree on the right ----
  \node[troot] (r) at (0,0) {\scriptsize root (split)};
  \node[tleaf] (tl) at (-2.6,-1.6) {\scriptsize leaf $0$};
  \node[tleaf] (tr) at (-0.9,-1.6) {\scriptsize leaf $0$};
  \node[tnode] (bl) at (0.9,-1.6) {\scriptsize leaf $1$};
  \node[tnode] (br) at (2.6,-1.6) {\scriptsize split};
  \node[tleaf] (b1) at (1.6,-3) {\scriptsize $0$};
  \node[tleaf] (b2) at (2.4,-3) {\scriptsize $0$};
  \node[tleaf] (b3) at (3.2,-3) {\scriptsize $0$};
  \node[tleaf] (b4) at (4.0,-3) {\scriptsize $0$};
  \draw[tedge] (r) -- (tl); \draw[tedge] (r) -- (tr);
  \draw[tedge] (r) -- (bl); \draw[tedge] (r) -- (br);
  \draw[tedge] (br) -- (b1); \draw[tedge] (br) -- (b2);
  \draw[tedge] (br) -- (b3); \draw[tedge] (br) -- (b4);
\end{tikzpicture}
\caption{The same grid, two representations. The $4\times4$ grid on the left is split (dashed) into
four $2\times2$ quadrants; three are uniform (leaf $0$, leaf $0$, leaf $1$) and the bottom-right
still mixes $0$/$1$, so it recurses one more level. The quad tree on the right is the
\emph{reconstruction}: a node is a leaf exactly when its square is uniform, internal otherwise --
LC 427.}\label{fig:c25-quad}
\end{figure}

\begin{definition}{Region quad tree}{c25-quaddef}
A \term{quad tree} over an $n\times n$ grid is a tree whose root covers the whole square. A node
covering a square $S$ is a \term{leaf} if every cell of $S$ holds the same value (that value is
stored on the leaf); otherwise it is \term{internal} and has exactly four children covering the
top-left, top-right, bottom-left and bottom-right halves of $S$.
\end{definition}

\subsection{Variations and Twists}

\begin{itemize}[nosep]
  \item \textbf{Sorted vs labelled degrees}: \thmref{c25-degreal} is about the labelled model (node
        $i$ has degree $d_i$). If only the \emph{multiset} of degrees matters, divide the count by
        automorphisms -- still a Pr\"ufer argument (\chref{c32}). \hfill\lv{4}--\lv{5}
  \item \textbf{Bounded-degree trees}: constrain which $d_i$ are allowed (e.g.\ binary: all $d_i\le3$
        for a rooted binary tree); realisability then needs a matching argument, not just the sum. \hfill\lv{5}
  \item \textbf{Prefer a specific tree}: to \emph{rebuild $T$ itself}, the degree array is not enough --
        you need a full Pr\"ufer sequence (the order matters), \chref{c32}. \hfill\lv{4}
  \item \textbf{Quad tree, faster uniformity test}: replace the $O(\text{area})$ scan with a 2-D prefix
        sum of the values so ``is $S$ uniform?'' is $O(1)$; total build $\bigO{n^{2}}$. \hfill\lv{4}
  \item \textbf{k-d tree / octree}: same partition-tree idea with $k$ dimensions or splits on
        a data point instead of uniformity -- construction-by-recursion, LC 427's template. \hfill\lv{3}--\lv{4}
\end{itemize}

\subsection{Common Pitfalls}\label{sec:c25-degree-pitfalls}

\begin{warningbox}[Trap 1 -- accepting a sequence with a zero or a bad sum]
$d_i=0$ with $n\ge2$ is impossible in a connected graph (an isolated node), and $\sum d_i\ne2(n-1)$
is fatal -- yet both sneak past code that only ``starts peeling''. Test positivity \emph{and} the sum
\emph{before} constructing (see the two early \code{return \{\}} lines of \coderef{c25-degbld}).
\end{warningbox}
\begin{lstlisting}[caption={\textcolor{warningred}{WRONG}: peel without validating}, language=C++, basicstyle=\ttfamily\small]
for (int i=0;i<n;i++) if (deg[i]==1) leaf.push(i);   // sum may be != 2(n-1): heap can empty mid-run
\end{lstlisting}
\begin{lstlisting}[caption={\textcolor{successgreen}{RIGHT}: check the certificate first}, language=C++, basicstyle=\ttfamily\small]
long long s=0; for(int d:deg){ if(d<1) return {}; s+=d; }
if (s != 2LL*(n-1)) return {};                        // reject before any peel
\end{lstlisting}

\begin{warningbox}[Trap 2 -- assuming the constructed tree is \emph{the} tree]
A degree sequence fixes only the \emph{degrees}, not the wiring. Unless a Pr\"ufer sequence's
\emph{order} is given, you recover \emph{one of} the $\prod$-many realisations -- fine for
``build any'', wrong for ``recover exactly $T$''. Do not claim the answer is unique when
\propref{c25-prufdeg} returns a count $>1$.
\end{warningbox}

\begin{warningbox}[Trap 3 -- quad tree: splitting uniform squares]
The base case is \emph{uniform value}, not \emph{size~$1$}. A $2\times2$ block of all $1$s is a single
leaf; recursing past a uniform square spawns empty children and corrupts \code{isLeaf}. Test
uniformity \emph{first}, and only split when it fails (\coderef{c25-quad}).
\end{warningbox}

\subsection{Problem Lab}

The anchor, LC 427, is the purest ``data $\to$ tree'' reconstruction in the chapter: the edges are
never mentioned -- they are \emph{implied} by the recursive partition of a grid.

% ---------------------------------------------------------------------------
\solvedproblem{c25-quad}{Construct Quad Tree}{LeetCode 427}
  {https://leetcode.com/problems/construct-quad-tree/}{LC Medium L3}
  {divide and conquer, recursive construction, uniformity test, tree building from grid}

\begin{problemstatement}[Construct Quad Tree]
Given an $n\times n$ grid \code{grid} of $0$s and $1$s, build its \textbf{quad tree} and return the
root. A \textbf{leaf} has \code{isLeaf=true} and \code{val} equal to the (single) value in its
region. An \textbf{internal} node has \code{isLeaf=false}, arbitrary \code{val}, and four children
\code{topLeft}, \code{topRight}, \code{bottomLeft}, \code{bottomRight} covering the four quadrants.
$n$ is a power of~$2$, $1\le n\le 6$.
\end{problemstatement}

\begin{attemptgate}[Attempt it yourself first -- \lv{3}, give yourself 25 minutes]
Fair first attempt: recursion on a square. Base case when the whole square is one value.
\peek{How do you decide ``the square $(r_0,c_0)\!\to\!(r_1,c_1)$ is uniform''? And when it
\emph{isn't}, into how many sub-squares does it break, and what are their corners?}
\nudge{Scan the square; if all cells equal $\mathit{val}$, return \code{new Node(val, true)}.
Otherwise split at the mid row/col into four rectangles -- top-left $[r_0,rm)\times[c_0,cm)$, etc.\
-- and recurse; return \code{new Node(val, false, tl, tr, bl, br)}.}
\fullhint{\code{build(r0,c0,r1,c1)}: let $v=\mathit{grid}[r_0][c_0]$; if the region is all $v$
return a leaf; else $rm=(r_0{+}r_1)/2$, $cm=(c_0{+}c_1)/2$ and recurse on the four half-open
quadrants, threading them as the four children. Top call \code{build(0,0,n,n)}. Full build and
proof below.}
\end{attemptgate}

\solstep{Understanding the Problem}
The tree's shape is completely determined by the grid: a node is a leaf iff its region is uniform.
The four children are ordered, so this is a \emph{fixed} 4-ary recursion -- the only real work is
(a) testing uniformity and (b) getting the quadrant boundaries right (half-open rectangles so the
four pieces tile the parent exactly). Unlike \secref{c25-parent} there are no edges to read -- the
edges are \emph{created} by the partition.

\solstep{Key Observations}
\begin{observation}
$n$ is a power of two, so halving always lands on integer boundaries and bottoms out at $1\times1$
(a trivially uniform square), guaranteeing the recursion terminates.
\end{observation}
\begin{observation}
``Uniform'' is a property of a region, checked independently at every node; the naive scan is
$O(\text{area})$. Because LC 427 caps $n\le6$, the naive total $O(n^2\log n)$ is trivially fine --
but a 2-D prefix sum makes each test $O(1)$ and the whole build $O(n^2)$ for large grids.
\end{observation}
\begin{observation}
For an internal node the \code{val} is unspecified by the problem, so any consistent choice works;
the standard is the top-left cell's value -- it must not be mistaken for a leaf decision.
\end{observation}

\solstep{Think It Through With Me}
\begin{thinkbox}
\deadend{Build the tree by first converting the grid into an explicit edge list.}
{There \emph{are} no edges in the input -- you would invent a graph, then re-derive the very
partition you were trying to encode. Quad trees are built \emph{top-down by recursion}, not from pairs.}
\deadend{Split before testing uniformity.}{Then a solid $4\times4$ block of $1$s becomes $16$ tiny
leaves and four levels of pointless children -- you produce a non-minimal quad tree and fail the
\code{isLeaf} checks. Uniformity is the \emph{base case}; test it first.}
\aha{A node is a leaf \emph{iff} its square is uniform. So: test, then either return a leaf or split
into the four half-open quadrants and recurse -- the recursion \emph{is} the construction.}
\end{thinkbox}

\solstep{Solution}
Define $\mathit{build}(r_0,c_0,r_1,c_1)$ over the half-open square $[r_0,r_1)\times[c_0,c_1)$. Read
$v=\mathit{grid}[r_0][c_0]$; if every cell in the square equals $v$, return a leaf \code{Node(v,true)}.
Otherwise compute midpoints $rm,cm$ and return an internal \code{Node(v,false,\ldots)} built from
the four recursive calls on the top-left $[r_0,rm)\times[c_0,cm)$, top-right $[r_0,rm)\times[cm,c_1)$,
bottom-left $[rm,r_1)\times[c_0,cm)$ and bottom-right $[rm,r_1)\times[cm,c_1)$ quadrants. The root
is $\mathit{build}(0,0,n,n)$.

\solstep{Proof of Correctness}
\begin{proofbox}
\emph{Termination}: each recursive call halves the side length, reaching $1\times1$ (uniform,
hence a leaf), so the depth is $\log_2 n$ and every branch terminates. \emph{Leaf correctness}:
$\mathit{build}$ returns \code{isLeaf=true} only after scanning the whole region and confirming all
cells equal $v$, matching the definition. \emph{Internal correctness}: when the uniformity scan
fails, the four half-open quadrants are disjoint and their union is exactly the parent square, so the
four children cover the parent's region with no gaps or overlaps, each built recursively to
represent its sub-region. By induction on region size, every returned node faithfully represents its
square, and the root represents the full grid -- so the constructed quad tree is the required one. \end{proofbox}

\solstep{Implementation}
\begin{lstlisting}[caption={Construct Quad Tree -- recursive partition of a uniform test (LeetCode 427)}, label={lst:c25-quad}]
#include <bits/stdc++.h>
using namespace std;

// Definition for a quad-tree node.
class Node {
public:
    bool val;
    bool isLeaf;
    Node* topLeft;
    Node* topRight;
    Node* bottomLeft;
    Node* bottomRight;
    Node() : val(false), isLeaf(false),
             topLeft(nullptr), topRight(nullptr),
             bottomLeft(nullptr), bottomRight(nullptr) {}
    Node(bool _val, bool _isLeaf) : val(_val), isLeaf(_isLeaf),
             topLeft(nullptr), topRight(nullptr),
             bottomLeft(nullptr), bottomRight(nullptr) {}
    Node(bool _val, bool _isLeaf, Node* a, Node* b, Node* c, Node* d)
        : val(_val), isLeaf(_isLeaf), topLeft(a), topRight(b), bottomLeft(c), bottomRight(d) {}
};

class Solution {
    vector<vector<int>> g;
    bool uniform(int r0, int c0, int r1, int c1, int v) {
        for (int i = r0; i < r1; ++i)
            for (int j = c0; j < c1; ++j)
                if (g[i][j] != v) return false;
        return true;
    }
    Node* build(int r0, int c0, int r1, int c1) {
        int v = g[r0][c0];
        if (uniform(r0, c0, r1, c1, v)) return new Node(v == 1, true);   // leaf (base case)
        int rm = (r0 + r1) / 2, cm = (c0 + c1) / 2;                     // split into 4
        return new Node(v == 1, false,
            build(r0, c0, rm, cm), build(r0, cm, rm, c1),              // top-left, top-right
            build(rm, c0, r1, cm), build(rm, cm, r1, c1));             // bottom-left, bottom-right
    }
public:
    Node* construct(vector<vector<int>>& grid) {
        g = grid;
        int n = (int)grid.size();
        return build(0, 0, n, n);
    }
};
\end{lstlisting}

\solstep{Trace}
\begin{tracebox}[The \figref{c25-quad} grid]
\figref{c25-quad}: bottom-left $2\times2$ is all~$1$; top-left and top-right $2\times2$ are all~$0$;
bottom-right is mixed. \code{build(0,0,4,4)} scans the whole $4\times4$ -- mixed $\Rightarrow$ internal.
Four calls: \code{build(0,0,2,2)} uniform $0\Rightarrow$ leaf; \code{build(0,2,2,4)} uniform $0\Rightarrow$
leaf; \code{build(2,0,4,2)} uniform $1\Rightarrow$ leaf; \code{build(2,2,4,4)} mixed $\Rightarrow$ split
into four $1\times1$ leaves $(1,0,0,0)$. The returned tree has root with $3$ leaf children and one
internal child with $4$ leaves -- exactly \figref{c25-quad}. \checkmark
\end{tracebox}

\solstep{Complexity}
\begin{complexitybox}
Time $O(n^2\log n)$ with the naive scan: level $0$ scans $n^2$, the $4$ children together scan
$4\cdot(n/2)^2=n^2$, and so on for $\log n$ levels -- $O(n^2)$ per level. Space $O(n^2)$ nodes in the
worst checkerboard (every $1\times1$ is its own leaf, $O(n^2)$ nodes). With $n\le6$ this is free; for
large grids, a 2-D prefix sum makes \code{uniform} $O(1)$ and the build $O(n^2)$. Verdict: \ac.
\end{complexitybox}

\solstep{Pitfalls and Variants}
\begin{warningbox}
Three breakers: (i) splitting a uniform region (Trap~3) -- test \emph{first}. (ii) Off-by-one
quadrants: use half-open $[r_0,r_1)$ and midpoints so the four children tile exactly; closed
boundaries overlap and one cell gets two parents. (iii) Forgetting \code{isLeaf} on internal nodes
(\code{new Node(v,false,\dots)}) -- the judge compares structure, so a mis-flagged leaf fails.
\emph{Variant:} LC 427's sibling is the ``intersection of two quad trees'' and ``quad-tree
compression'' -- same recursion, extra combine step.
\end{warningbox}

\begin{insightbox}[Key takeaways]
\begin{enumerate}[nosep, leftmargin=1.6em]
  \item \textbf{Degrees are a certificate.} One equation, $\sum d_i=2(n-1)$, decides existence
        outright (\thmref{c25-degreal}) -- far cleaner than the general-graph case.
  \item \textbf{The peel \emph{is} Pr\"ufer decoding.} Leaf-peel construction and Pr\"ufer counting
        are the same mechanism seen constructively vs. enumeratively (\chref{c32}).
  \item \textbf{Recursive partition is reconstruction too.} A quad tree turns grid data into a tree
        with no edges ever named -- construction from aggregate data, one shape per data value.
\end{enumerate}
\end{insightbox}

\begin{hardeningbox}
\begin{itemize}[nosep, leftmargin=1.4em]
  \item \textbf{O(1) uniformity}: build a 2-D prefix sum \code{pre} of the grid; a region is uniform
        iff its cell-sum equals $0$ or equals its area -- turns \probref{c25-quad} into $O(n^2)$.
  \item \textbf{Recover \emph{exactly} $T$}: replace \coderef{c25-degbld}'s chosen Pr\"ufer order with
        the given code, then decode -- the standard Prüfer$\to$tree decoder (\chref{c32}).
  \item \textbf{Existential edge count}: among the $180$ realisations of $T$'s degrees, ask how many
        have node $2$ adjacent to node $7$; a Pr\"ufer-window counting exercise, \chref{c32}.
\end{itemize}
\end{hardeningbox}

\subsection{Practice Problems}

\begin{stuckbox}
\begin{itemize}[nosep, leftmargin=1.4em]
  \item Did you validate positivity \emph{and} the $\sum=2(n-1)$ equation before peeling?
        (Trap~1, \thmref{c25-degreal}.)
  \item Did you resist claiming uniqueness when \propref{c25-prufdeg} gives a count $>1$? (Trap~2.)
  \item In the quad tree, did you test uniformity \emph{before} splitting? (Trap~3, \probref{c25-quad}.)
  \item Are your quadrant rectangles half-open so the four children tile the parent? (\probref{c25-quad}.)
\end{itemize}
\end{stuckbox}

\begin{practiceset}
  \practice{Course Schedule II}{LeetCode 210}
    {https://leetcode.com/problems/course-schedule-ii/}{LC Medium L3}
    {Given prerequisite \emph{pairs}, return a valid ordering of all courses, or \code{[]}. This is the
    \textbf{in-degree peel} of \thmref{c25-degreal} on a \emph{directed acyclic} graph: repeatedly emit
    every node of current in-degree~$0$ (a ``source'', the DAG analogue of a degree-$1$ leaf), delete
    its outgoing edges, and repeat (Kahn's algorithm). Recognise the shared mechanic -- peel by degree,
    reconstruct an order/structure from aggregate counts -- and note the \emph{cycle detector}: if the
    peel finishes with nodes still having positive in-degree, the input is not a DAG (the ``realisability
    fails'' moment). Rebuild \code{adj} from the pairs first, exactly \secref{c25-parent}.}
\end{practiceset}

\begin{recallbox}[10-second refresher -- 6.2 Degrees, Sequences \& Codes]
A positive list is a tree degree sequence \emph{iff} $\sum d_i=2(n-1)$ (\thmref{c25-degreal}); build a
realisation by peeling the smallest leaf onto a degree-$>1$ node (\coderef{c25-degbld}) -- the same
move as Pr\"ufer decoding, where node $v$ appears $\degop(v)-1$ times (\defref{c25-prufer}). The count
of realising labelled trees is $(n-2)!/\prod(d_i-1)!$ (\propref{c25-prufdeg}); full bijection and
Cayley are \chref{c32}. A quad tree reconstructs a grid by uniform-then-split recursion (\probref{c25-quad});
a topological order reconstructs a DAG by the in-degree peel (LC 210).
\selftest{
  \item Decide realisability of $[3,2,2,1,1,1]$ for $n=6$: yes or no, and by which test?
  \item Run one peel step on $T$'s degrees -- which node do you pin first and to what?
  \item How many labelled trees share a degree array with $\degop-1$ multiset $\{1,2,2,1\}$?
  \item In LC 427, why must uniformity be tested \emph{before} you split a square?
  \item Explain how Kahn's algorithm is the directed in-degree sibling of the leaf-peel.
}
\end{recallbox}

% !TEX root = ../../tree_notes.tex
% =====================================================================
%  6.3  Is It a Tree? Identification + Union-Find Pre-Look  sec:c25-identify
%  Exports: def:c25-comp  prop:c25-uf  cor:c25-istree  lst:c25-uf  ex:c25-id-run
%           prob:c25-rc1  lst:c25-rc1  (practice: LC 547)
% =====================================================================
\section{Is It a Tree? Identification and a Union-Find Pre-Look}\label{sec:c25-identify}

\begin{previouslybox}
\begin{enumerate}[nosep, leftmargin=1.6em]
  \item State the tree certificate and explain why checking only \emph{two} of its three legs
        suffices. (\thmref{c25-cert})
  \item In \probref{c25-rc2}, what did a union-find \emph{rejection} (an edge whose endpoints were
        already connected) tell you? (\probref{c25-rc2})
  \item A forest with $n$ nodes and $k$ components has how many edges? Guess, then prove it in
        this section. (\thmref{c25-degsum})
  \item Why does banning the \emph{later} of two candidate edges and testing for a surviving cycle
        pick the right redundant edge? (\probref{c25-rc2})
\end{enumerate}
\end{previouslybox}

\begin{topicintro}
Sections~6.1 and~6.2 \emph{built} a tree from trustworthy-ish data. This section is the decision
problem \emph{before} you build: \textbf{given an edge list, is it a tree at all -- or is it a
forest, or does it hide a cycle, and if there is a cycle, which edge is redundant?} The
union-find you touched in \probref{c25-rc2} becomes the star: process edges left to right, merging
components, and the very first edge whose endpoints are \emph{already} in one set is a
\term{cycle-closing} edge. That single observation \emph{proves} the tree certificate operationally,
\emph{counts components} for free (the ``number of provinces'' problem), and \emph{identifies and
deletes} the one surplus edge (LC 684 -- the undirected little sibling of LC 685). This is a
\textbf{pre-look}: enough union-find to certify trees. The \emph{full} data-structure -- union by
rank and path compression analysis -- lives in \chref{c28} (its use in Kruskal is Ch.~31,
\chref{c17}); we borrow only the
``are these two already connected?'' query here.
\textbf{Prerequisites.} \thmref{c25-cert} (the certificate), \probref{c25-rc2} (the union-find
cycle alarm). This section formalises \emph{why} that alarm is exactly acyclicity and turns it into
a decision procedure.

\textbf{After this section you can}
\begin{itemize}[nosep]
  \item decide ``is this a tree?'' in $\bigO{n}$ with the $n-1$-edge count plus one union-find pass; \hfill\lv{2}--\lv{3}
  \item classify any edge set as tree / forest / unicyclic / multi-cyclic by counting rejections;  \hfill\lv{3}
  \item count connected components (provinces) with the same machinery;                              \hfill\lv{2}--\lv{3}
  \item find and delete the last redundant edge that keeps a graph a tree (LC 684).                 \hfill\lv{3}
\end{itemize}
\end{topicintro}

\begin{recognitionbox}[Pattern card -- certify, classify, or repair a suspicious edge set]
\reckey{you are given \code{edges} on $n$ nodes and the question is about \emph{validity or shape},
not a subtree computation: ``is it a tree?'', ``how many connected pieces?'', ``which single edge
makes it cyclic?'', ``delete one edge to get a tree''.}
\recuse{
  \item ``return the redundant connection'' -- the edge closing a cycle, last-in-input wins (\probref{c25-rc1});
  \item ``number of provinces / friend circles / islands'' -- count components (\secref{c25-identify});
  \item ``can all pairs be connected / is this network a tree?'' -- $n-1$ edges + connectivity;
  \item the directed / two-parent version -- \probref{c25-rc2} already used this pre-look.
}
\recnotuse{
  \item you only need to build \code{adj} for a trusted tree -- \secref{c25-parent};
  \item you need the \emph{minimum spanning tree} weight, not the yes/no shape -- \chref{c17};
  \item connectivity under edge \emph{insertions and deletions} online -- that needs more than this pre-look.
}
One-line test: \emph{is the question ``does this edge set form one acyclic whole?'' Merge
with union-find and watch for a rejection -- this section.}
\end{recognitionbox}

\subsection{Intuition}

A tree is \emph{maximally acyclic} and \emph{minimally connected}: add any edge and you create a
cycle; delete any edge and you disconnect. That dual tightness is the certificate of
\thmref{c25-cert}, and union-find is the machine that \emph{enforces} it edge by edge. Feed edges in
order, each time asking ``are these two endpoints already reachable from each other?'' -- if yes,
this edge is \emph{redundant} (it closes a cycle); if no, merge the two pieces. After all edges: if
nothing was rejected and you used exactly $n-1$ edges, you watched $n$ separate nodes coalesce into
one -- a tree. A single rejection is a cycle; a leftover set of pieces is a forest.

\begin{figure}[H]
\centering
\begin{tikzpicture}[diagram, scale=0.82, transform shape]
  % ---- tree ----
  \begin{scope}[xshift=-4.6cm]
    \node[annot] at (1.3,0.7) {\small tree ($e=n-1$, 0 rejects)};
    \node[troot] (a) at (0,0) {$1$};
    \node[tnode] (b) at (-1,-1.2) {$2$};
    \node[tnode] (c) at (1,-1.2) {$3$};
    \node[tleaf] (d) at (-1,-2.4) {$4$};
    \draw[tedgehl] (a)--(b); \draw[tedgehl] (a)--(c); \draw[tedgehl] (b)--(d);
  \end{scope}
  % ---- forest ----
  \begin{scope}[xshift=0cm]
    \node[annot] at (1,0.7) {\small forest ($e=n-2$, 2 pieces)};
    \node[tnode] (f1) at (0,0) {$1$};
    \node[tnode] (f2) at (1,0) {$2$};
    \node[tnode] (f3) at (2.4,0) {$3$};
    \draw[tedgehl] (f1)--(f2);
    \node[annot, font=\scriptsize] at (2.4,-0.9) {isolated};
  \end{scope}
  % ---- unicyclic ----
  \begin{scope}[xshift=4.4cm]
    \node[annot] at (0.6,0.7) {\small unicyclic ($e=n$, 1 reject)};
    \node[tnode] (u1) at (0,0) {$1$};
    \node[tnode] (u2) at (1.2,0) {$2$};
    \node[tnode] (u3) at (0.6,-1.2) {$3$};
    \draw[tedgehl] (u1)--(u2); \draw[tedgehl] (u2)--(u3);
    \draw[tedge, draw=warningred, line width=1.1pt] (u3)--(u1);
    \node[annot, font=\scriptsize, text=warningred] at (0.14,-1.05) {close};
  \end{scope}
\end{tikzpicture}
\caption{The three shapes union-find separates. A \textbf{tree}: $n-1$ edges, every merge accepted.
A \textbf{forest}: acyclic but $<n-1$ edges, several surviving sets. A \textbf{unicyclic} graph: $n$
edges, the red edge finds its endpoints already joined and is \emph{rejected} -- the redundant
connection (\probref{c25-rc1}). The number of rejects is the cyclomatic number $e-(n-k)$.}\label{fig:c25-shapes}
\end{figure}

\subsection{Definitions and Notation}

\begin{definition}{Forest, component, spanning tree}{c25-comp}
A \term{forest} is an acyclic graph. Its maximal connected subgraphs are \term{components}; a
connected forest (one component) is a tree. A graph on $n$ nodes with $k$ components that is a
forest has exactly $n-k$ edges. A \term{spanning tree} of a connected graph uses all $n$ nodes and
exactly $n-1$ of its edges.
\end{definition}

\begin{proposition}{Union-find pre-look at acyclicity}{c25-uf}
Process a graph's $m$ edges one at a time, initially $n$ singleton sets, merging (unioning) the two
endpoint sets for each \emph{accepted} edge and \emph{rejecting} any edge whose endpoints already
share a set. Then:
\begin{enumerate}[nosep, leftmargin=1.6em]
  \item an edge is rejected \textbf{iff} it closes a cycle among the edges seen so far;
  \item the number of accepted edges is exactly $n-k$, where $k$ is the number of connected
        components of the whole graph;
  \item the number of rejections is $m-(n-k)$ -- the graph is a forest \textbf{iff} this is $0$.
\end{enumerate}
\end{proposition}
\begin{proofbox}
(1) At the moment $[u,v]$ is examined, the accepted edges define reachability; $u,v$ already share a
set precisely when a path of accepted edges connects them, so adding $[u,v]$ closes a cycle.
Conversely any cycle's last-processed edge sees its two neighbours already joined. (2) Each accepted
union lowers the set count by exactly one, from $n$ down; the surviving sets are the connected
components of the graph (accepted edges and rejected edges connect the same pairs, since a rejected
edge joined already-connected nodes), so exactly $k$ sets remain and $n-k$ unions were accepted.
(3) Rejections $=m-\text{accepted}=m-(n-k)$. A forest has no cycles, hence no rejections; a graph
with no rejections never closed a cycle, hence is acyclic -- a forest. \end{proofbox}

\begin{corollary}{The operational tree test}{c25-istree}
A graph on $n$ nodes is a tree \textbf{iff} it has $n-1$ edges and union-find rejects none, and
\textbf{iff} it has $n-1$ edges and union-find merges everything into a single set. This is
\thmref{c25-cert} turned into a one-pass algorithm: the count is leg~(1), the rejection-free scan
is leg~(3), and either plus the count forces leg~(2).
\end{corollary}
\begin{insightbox}
You get \emph{three} answers from one scan for free: reject-count $>0$ means ``not a tree, has a
cycle''; reject-count $=0$ but sets${}_\text{final}>1$ means ``a forest''; reject-count $=0$ and
exactly one final set with $n-1$ edges means ``a tree''. LC 684 reads off the \emph{last} rejected
edge; LC 547 reads off the \emph{final set count}.
\end{insightbox}

\subsection{The algorithm: one union-find pass to certify and classify}

\begin{lstlisting}[caption={Union-find pre-look: decide tree / forest / cyclic, and count components}, label={lst:c25-uf}]
#include <bits/stdc++.h>
using namespace std;

struct DSU {
    vector<int> p;
    DSU(int n) : p(n) { iota(p.begin(), p.end(), 0); }        // p[i]=i: n singleton sets
    int find(int x) { return p[x] == x ? x : p[x] = find(p[x]); }   // path compression
    bool uni(int a, int b) {                                  // false => already connected
        a = find(a); b = find(b);
        if (a == b) return false;                             // rejected edge: closes a cycle
        p[b] = a; return true;                                // accepted: merges two sets
    }
};

// Is the graph on nodes 0..n-1 exactly ONE tree?
bool isTree(int n, const vector<pair<int,int>>& edges) {
    if ((int)edges.size() != n - 1) return false;            // leg (1): the edge count
    DSU d(n);
    for (auto& e : edges)
        if (!d.uni(e.first, e.second)) return false;         // leg (3): a reject = a cycle
    return true;                                             // n-1 edges + acyclic => connected
}

// Number of connected components (provinces).
int countComponents(int n, const vector<pair<int,int>>& edges) {
    DSU d(n);
    int comps = n;
    for (auto& e : edges) if (d.uni(e.first, e.second)) --comps;  // each merge removes one piece
    return comps;
}
\end{lstlisting}
\begin{complexitybox}
Time $\bigO{(n+m)\,\alpha(n)}$ -- every \code{find}/\code{uni} is inverse-Ackermann with path
compression ($\alpha\le4$ in practice), so essentially linear. Memory $\bigO{n}$. This one-pass
approach is far cheaper than the alternative DFS colouring for cycle detection, and \emph{both}
answers (tree? how many components?) come from one pass. Full DSU complexity analysis (why path
compression alone is already near-linear) is \chref{c28}.
\end{complexitybox}

\subsection{Worked Example}

\begin{example}{Classifying $T$ and $T$ plus one bad edge}{c25-id-run}
$T$ on $n=8$ has $7=n-1$ edges. \code{isTree}: the count passes, and processing the seven edges in
\{1-2,1-3,2-4,2-5,3-6,3-7,7-8\} merges $8\!\to\!1$ set with no rejection -- \textbf{a tree}
\checkmark, \code{countComponents}$=1$. Now append the bad edge $[4,5]$: $8$ edges for $8$ nodes
fails the count leg instantly, and if you ran the scan anyway, edges up to $7$ build one set, then
$[4,5]$ finds $4,5$ already joined $\Rightarrow$ \textbf{reject} (one cycle, the
$2\text{--}4\text{--}5$ triangle). The \emph{last} rejected edge is exactly $[4,5]$ -- the edge LC
684 would return.
\end{example}
\begin{tracebox}[Edge-by-edge on LC 684's Example~2]
\code{edges=[[1,2],[2,3],[3,4],[1,4],[1,5]]}, nodes $1..5$. Sets start $\{1\}\{2\}\{3\}\{4\}\{5\}$.
$[1,2]$: merge $\{1,2\}$. $[2,3]$: merge $\{1,2,3\}$. $[3,4]$: merge $\{1,2,3,4\}$. $[1,4]$:
$\mathit{find}(1)=\mathit{find}(4)$ -- \textbf{already one set} $\Rightarrow$ reject, return $[1,4]$. The scan stops
at the first rejection; $[1,4]$ is the redundant connection, matching the official answer. \checkmark
\end{tracebox}

\subsection{Variations and Twists}

\begin{itemize}[nosep]
  \item \textbf{DFS cycle test}: root the graph and watch for a back-edge to a non-parent ancestor --
        equivalent to union-find on undirected graphs, $\bigO{n+m}$; union-find is less code. \hfill\lv{3}
  \item \textbf{Directed version}: acyclicity now means a \emph{topological order} exists (LC 207),
        and the peel is in-degree based (\secref{c25-degree}'s practice) -- union-find does NOT apply
        to directed cycles. \hfill\lv{3}--\lv{4}
  \item \textbf{Multiple surplus edges}: to delete $k$ edges and leave a forest, \emph{count} rejections;
        a graph is a $k$-cyclomatic forest iff rejections $=k$. \hfill\lv{4}
  \item \textbf{Dynamic connectivity}: edges arriving over time, ask ``connected yet?'' -- union-find
        answers monotone insertions directly (this pre-look already does); deletions need more (\chref{c28}). \hfill\lv{4}
  \item \textbf{Bipartite / 2-colour check}: union-find with parity answers ``can nodes be split into
        two sets across every edge'' -- the same merge mechanic with an extra flag (\chref{c28}). \hfill\lv{4}
\end{itemize}

\subsection{Common Pitfalls}\label{sec:c25-identify-pitfalls}

\begin{warningbox}[Trap 1 -- trusting the edge count alone]
``$n-1$ edges'' does \emph{not} mean ``tree'': $n$ nodes with $n-1$ edges could be a disconnected
piece plus a small cycle elsewhere (the two defects cancel in the count). Union-find rejects none
\emph{and} the count is $n-1$ -- both legs (\corref{c25-istree}); or check connectivity instead of
acyclicity.
\end{warningbox}
\begin{lstlisting}[caption={\textcolor{warningred}{WRONG}: count only -- a unicyclic disconnected graph slips past}, language=C++, basicstyle=\ttfamily\small]
return (edges.size() == n-1);        // 4 nodes, edges {1-2,2-3,3-1}: count 3? no. See the trap.
\end{lstlisting}
\begin{lstlisting}[caption={\textcolor{successgreen}{RIGHT}: count + rejection-free union-find}, language=C++, basicstyle=\ttfamily\small]
if ((int)edges.size() != n-1) return false;
DSU d(n);
for (auto& e : edges) if (!d.uni(e.first, e.second)) return false;   // a reject = cycle
return true;
\end{lstlisting}

\begin{warningbox}[Trap 2 -- union-find on a \emph{directed} cycle]
The rejection test detects \emph{undirected} cycles. On the directed graph of LC 207 the two
prerequisite edges $a\!\to\!b$ and $b\!\to\!a$ are a genuine directed cycle but union-find would merge
them silently the same as any pair. Detecting directed cycles needs Kahn's in-degree peel or DFS
colour -- not this tool (\secref{c25-degree}).
\end{warningbox}

\begin{warningbox}[Trap 3 -- forgetting path compression / recursion depth]
Without \code{p[x]=find(p[x])} the trees get tall; at $n\approx2\cdot10^{5}$ a \emph{recursive}
\code{find} on a degenerate chain can even overflow the stack. Keep compression (as in
\coderef{c25-uf}) or make \code{find} iterative; the union-by-rank upgrade is \chref{c28}.
\end{warningbox}

\subsection{Problem Lab}

LC 684 is the section distilled: one undirected extra edge, return the last redundant one -- the
\emph{first} rejection of a union-find run in input order \emph{is} the answer.

% ---------------------------------------------------------------------------
\solvedproblem{c25-rc1}{Redundant Connection}{LeetCode 684}
  {https://leetcode.com/problems/redundant-connection/}{LC Medium L3}
  {undirected graph, union-find cycle detection, spanning tree, edge deletion}

\begin{problemstatement}[Redundant Connection]
A \textbf{tree} on $n$ nodes (labels $1\dots n$) gained exactly \textbf{one} extra \emph{undirected}
edge, so \code{edges} now has $n$ edges (each \code{[u,v]} with $u\ne v$). Find the edge that can be
removed to make the graph a tree again; if several answers exist, return the one appearing
\textbf{last} in \code{edges}. $2\le n\le 1000$.
\end{problemstatement}

\begin{attemptgate}[Attempt it yourself first -- \lv{3}, give yourself 20 minutes]
Fair first attempt: a tree has $n-1$ edges; here there are $n$, so exactly one is on a cycle. Find a
cycle, return its latest edge.
\peek{Process edges left to right. When you are about to add $[u,v]$, what simple question tells you
this edge is the one closing the cycle?}
\nudge{Maintain connected components with union-find. The first edge whose two endpoints are
\emph{already connected} closes the only cycle. Because you scan in input order, the \emph{last}
such edge you could point at is naturally the one you return -- there is exactly one cycle here.}
\fullhint{Run union-find over $1..n$ (arrays sized $n+1$) in input order; return the first edge
$[u,v]$ with $\mathit{find}(u)=\mathit{find}(v)$. That edge is redundant and, since the input has a single cycle and
you scanned in order, it is the latest edge of that cycle. Full build and proof below.}
\end{attemptgate}

\solstep{Understanding the Problem}
The graph has $n$ nodes and $n$ edges: one more edge than a tree needs, so it holds \emph{exactly
one} cycle (cyclomatic number $m-(n-k)=n-(n-1)=1$ with $k=1$). Removing any cycle edge yields a tree;
the tie-break wants the \emph{last} cycle edge in input order. A rooted/two-parent complication is
absent -- edges are undirected, so this is \probref{c25-rc2} with Case~2 only.

\solstep{Key Observations}
\begin{observation}
Scanning edges in order and merging with union-find, the \emph{first} edge whose endpoints are
already connected closes the cycle -- that is the unique redundant edge of a unicyclic graph, and it
appears \emph{last} among that cycle's edges by construction of the scan.
\end{observation}
\begin{observation}
You never need to find the cycle's \emph{members} -- only its closing edge. Union-find identifies the
closing edge directly, in $\bigO{n\,\alpha(n)}$, with no DFS and no parent bookkeeping.
\end{observation}
\begin{observation}
This is \corref{c25-istree} run just past the point of failure: a tree would reject nothing; the first
rejection here names the culprit.
\end{observation}

\solstep{Think It Through With Me}
\begin{thinkbox}
\deadend{DFS for the cycle, collect its edges, take the max-index one.}
{Workable but you must track a back-edge, reconstruct the cycle from parent pointers, then scan for
the largest input index -- three fiddly steps. The union-find scan \emph{is} the cycle-closing test,
in-order, so the ``last edge'' tie-break is automatic.}
\deadend{Check every edge by deleting it and re-testing $\mathit{isTree}$.}{$O(n)$ tests $\times$
$O(n)$ each $=O(n^{2})$; fine for $n\le1000$ but it re-derives connectivity $n$ times and still needs
careful tie-breaking. One pass beats it.}
\aha{Process edges in input order with union-find; return the \emph{first} rejected edge. Because
there is exactly one cycle, the closing edge is the last-in-order member of it -- the exact answer.}
\end{thinkbox}

\solstep{Solution}
Initialise union-find on $1..n$. Walk \code{edges} in order: for $[u,v]$, if $\mathit{find}(u)=\mathit{find}(v)$
return $[u,v]$ (the cycle-closing redundant edge); otherwise union them. Because a valid input
guarantees one extra edge, the loop always returns before finishing.

\solstep{Proof of Correctness}
\begin{proofbox}
The graph has $n$ edges on $n$ nodes, so it is connected with cyclomatic number $1$ -- a single
cycle $C$. Claim: the first edge $e^\ast$ the scan rejects is the member of $C$ appearing last in
\code{edges}. When $e^\ast=[u,v]$ is rejected, $u,v$ are already joined by accepted edges, forming a
$u$--$v$ path $P$ that avoids $e^\ast$, so $P\cup\{e^\ast\}$ is a cycle; the input has only one cycle
$C$, hence $P\cup\{e^\ast\}=C$ and $e^\ast\in C$. All earlier members of $C$ were accepted (they did
not reject, else the scan had stopped), and every edge processed after $e^\ast$ is not examined, so
$e^\ast$ is the last member of $C$ in input order. Removing $e^\ast$ breaks $C$ and leaves $n-1$
edges with no cycle -- a tree (\corref{c25-istree}) -- and it is the latest valid deletion, matching
the tie-break. \end{proofbox}

\solstep{Implementation}
\begin{lstlisting}[caption={Redundant Connection -- return the first edge union-find rejects (LeetCode 684)}, label={lst:c25-rc1}]
#include <bits/stdc++.h>
using namespace std;

class Solution {
    vector<int> p;
    int find(int x) { return p[x] == x ? x : p[x] = find(p[x]); }
public:
    vector<int> findRedundantConnection(vector<vector<int>>& edges) {
        int n = (int)edges.size();                 // n nodes, n edges (labels 1..n)
        p.resize(n + 1);
        iota(p.begin(), p.end(), 0);
        for (auto& e : edges) {
            int a = find(e[0]), b = find(e[1]);
            if (a == b) return e;                  // already connected => closes the cycle
            p[a] = b;                              // merge
        }
        return {};                                 // unreachable for valid input
    }
};
\end{lstlisting}

\solstep{Trace}
\begin{tracebox}[Both official examples]
\emph{Ex.~1} \code{[[1,2],[1,3],[2,3]]}: merge $1,2$; merge $1,3$ (set $\{1,2,3\}$); $[2,3]$:
$\mathit{find}(2)=\mathit{find}(3)$ $\Rightarrow$ return $[2,3]$. \checkmark\\
\emph{Ex.~2} \code{[[1,2],[2,3],[3,4],[1,4],[1,5]]}: merge $1,2,3,4$ into one set across the first
three edges; $[1,4]$ already connected $\Rightarrow$ return $[1,4]$ (matches expected). \checkmark
\end{tracebox}

\solstep{Complexity}
\begin{complexitybox}
Time $\bigO{n\,\alpha(n)}$ -- one in-order union-find pass ($\alpha$ inverse-Ackermann, $\le4$).
Space $\bigO{n}$ for the parent array. Beats the $O(n^{2})$ delete-and-retest brute force, and at
$n\le1000$ either passes. Verdict: \ac.
\end{complexitybox}

\solstep{Pitfalls and Variants}
\begin{warningbox}
(i) Size arrays $n+1$ for $1$-indexed labels (Trap~2 of \secref{c25-parent}). (ii) Return the
\emph{first} rejected edge, not the last rejected -- with a single cycle they coincide, but if you
kept scanning past the reject you might overwrite the answer; \code{return e} immediately.
(iii) Do not use this on the directed LC 685 without the in-degree pre-pass (\probref{c25-rc2}).
\emph{Variant:} LC 685 is this with a possible two-parent node; LC 687-style problems add a second
edge to delete -- re-run the scan after removing the found edge.
\end{warningbox}

\begin{insightbox}[Key takeaways]
\begin{enumerate}[nosep, leftmargin=1.6em]
  \item \textbf{A rejection \emph{is} a cycle.} Union-find turns ``is it a tree?'' into ``did I reject
        an edge?'' -- acyclicity for free, and count components on the same pass (\propref{c25-uf}).
  \item \textbf{Scan order = tie-break.} The first rejected edge is the last-in-order cycle member; the
        ``return latest'' rule needs no extra comparison.
  \item \textbf{684 $\subset$ 685.} This is \probref{c25-rc2}'s Case~2 -- undirected, one cycle, no
        two-parent node -- the certificate's legs checked one at a time.
\end{enumerate}
\end{insightbox}

\begin{hardeningbox}
\begin{itemize}[nosep, leftmargin=1.4em]
  \item \textbf{Two edges to delete}: run \probref{c25-rc1}, remove the found edge, re-run to find the
        next reject -- reject-count equals the cyclomatic number $m-(n-k)$.
  \item \textbf{Bigger inputs}: add union-by-rank to keep $\mathit{find}$ shallow (or make it iterative, Trap~3);
        full analysis \chref{c28}.
  \item \textbf{From tree to MSW}: if edges are weighted and you want the \emph{max}-weight edge to
        delete, sort by weight and run the same reject logic -- that is Kruskal in reverse (\chref{c17}).
\end{itemize}
\end{hardeningbox}

\subsection{Practice Problems}

\begin{stuckbox}
\begin{itemize}[nosep, leftmargin=1.4em]
  \item Did you check the $n-1$ count \emph{and} a rejection-free scan (not just one leg)?
        (\corref{c25-istree}, Trap~1.)
  \item Did you return on the \emph{first} rejection instead of scanning past it? (\probref{c25-rc1}.)
  \item Are your arrays sized for $1$-indexed nodes? (Trap~2 of \secref{c25-parent}.)
  \item For counting components, did you start the tally at $n$ and decrement on each accepted merge?
        (\coderef{c25-uf}, \probref{c25-rc1}.)
\end{itemize}
\end{stuckbox}

\begin{practiceset}
  \practice{Number of Provinces}{LeetCode 547}
    {https://leetcode.com/problems/number-of-provinces/}{LC Medium L2}
    {A $n\times n$ adjacency matrix \code{isConnected} of friend relationships; count the
    \textbf{provinces} (connected components). Rebuild the edge list from the matrix (read only
    $i<j$ to skip duplicates) and call \code{countComponents} from \coderef{c25-uf} -- start at $n$
    and decrement on each accepted union -- or DFS the matrix counting roots. This is exactly the
    ``reject-count $=0$ but sets${}>1$'' branch of \secref{c25-identify}: the answer is the number of
    surviving union-find sets (or DFS trees), and it needs \emph{no} edge-count leg since cycles are
    allowed. Recognise it as the component-counting half of this section.}
\end{practiceset}

\begin{recallbox}[10-second refresher -- 6.3 Is It a Tree? (Union-Find pre-look)]
Certify with \emph{two} certificate legs: exactly $n-1$ edges and a union-find pass that rejects
nothing (or merges to one set) (\corref{c25-istree}). A rejection is a cycle; rejections${}=m-(n-k)$
and $k$ is the final set count (\propref{c25-uf}). LC 684 returns the \emph{first} rejected edge in
input order -- the last cycle member (\probref{c25-rc1}); the directed sibling is LC 685
(\probref{c25-rc2}); components-province count is LC 547. Full DSU (rank, compression) is
\chref{c28}, and its use in Kruskal is Ch.~31 (\chref{c17}).
\selftest{
  \item Give an $n$-node edge set with $n-1$ edges that is \emph{not} a tree; which certificate leg fails?
  \item State the reject-count formula and use it to prove a graph with $0$ rejections is a forest.
  \item In LC 684, why is the \emph{first} rejected edge also the ``last appearing'' valid answer?
  \item Why must you NOT use union-find to detect a \emph{directed} cycle? (Name the right method.)
  \item Compute $\mathit{countComponents}$ for $T$ (8 nodes, 7 edges) and for $T$ with a self-loop added.
}
\end{recallbox}

% !TEX root = ../../tree_notes.tex
% =====================================================================
%  6.4  Recover the Structure by Queries / Interactive   sec:c25-query
%  Exports: def:c25-query  prop:c25-halve  thm:c25-logn  lst:c25-guess
%           ex:c25-query-run  prob:c25-guess  lst:c25-guess  (practice: LC 375)
% =====================================================================
\section{Recover the Structure by Queries (Interactive)}\label{sec:c25-query}

\begin{previouslybox}
\begin{enumerate}[nosep, leftmargin=1.6em]
  \item What does a single comparison in binary search do to the size of the candidate set? Why
        $\log_2 n$ comparisons in total? (\chref{c01}, \chref{c24} preview)
  \item In \secref{c25-parent}, what $\bigO{1}$ fact about the root made it readable from an array?
        (\propref{c25-root})
  \item How could a \emph{distance} or \emph{ancestor} answer about a tree be produced by the
        timestamp machinery? (\chref{c24}, \thmref{c24-anc})
  \item If a yes/no answer is truthful and splits the world in half, how many do you need to pin
        one of $N$ possibilities?
\end{enumerate}
\end{previouslybox}

\begin{topicintro}
Every topic so far \emph{handed} you the data. Here the data is hidden behind an \term{oracle}:
you may only \emph{ask questions}, and each answer is a small, costly bit of information you must
\emph{use adaptively} -- your next query depends on the previous answers. Reconstructing a hidden
object by queries is the same skill as the parent/degree/identification problems of
\secref{c25-parent}--\secref{c25-identify}, turned inside out: instead of reading a full array you
\emph{elicit} it, one strategic question at a time. The purest instance -- and our worked anchor
LC 374 -- is \term{guess the number}: pick from $[1,n]$, an oracle replies ``higher / lower /
correct'', and the optimal strategy is textbook binary search, which provably uses the minimum
possible number of questions (\thmref{c25-logn}). Two ideas carry the whole section: a truthful
answer should \emph{halve} the candidate space, and the number of candidates you can distinguish is
$2^{\text{\#answers}}$, so you cannot beat $\lceil\log_2\rceil$. Where the \emph{cost} of a wrong
guess matters (LC 375) the halving point shifts from the middle to a min-max choice -- an interval
DP -- and where the hidden object is a \emph{tree} rather than a number, the same adaptive idea
becomes group-testing on \term{distance}/\term{ancestor} queries answered by the Chapter~5 oracle
\chref{c24}. This is construction from data where you get to \emph{choose} the data.
\textbf{Prerequisites.} Binary-search intuition; the Chapter~5 distance/ancestor oracle
(\chref{c24}) used as a black box. This section adds the \emph{query model, the
information-theoretic lower bound, and adaptive vs.\ cost-aware strategies}.

\textbf{After this section you can}
\begin{itemize}[nosep]
  \item model a problem as an adaptive query protocol and identify the candidate space;        \hfill\lv{2}
  \item apply the halving principle and prove $\lceil\log_2 N\rceil$ questions are necessary;   \hfill\lv{2}--\lv{3}
  \item implement optimal search against a three-way (higher/lower/equal) oracle (LC 374);      \hfill\lv{2}
  \item recognise when the optimal split is \emph{not} the middle (worst-case cost, LC 375).    \hfill\lv{3}--\lv{4}
\end{itemize}
\end{topicintro}

\begin{recognitionbox}[Pattern card -- reconstruct by asking the oracle]
\reckey{you do not see the hidden value/object; instead you are given a \emph{feedback function}
(a \code{guess()}, an API that answers ``higher/lower'', a distance/ancestor query, a yes/no test)
and must find the answer with as few \emph{queries} as possible, adapting each to the last.}
\recuse{
  \item ``guess the number / word / root'' with a comparison oracle -- binary search (\probref{c25-guess});
  \item ``find $x$ using $f(x)$ queries, minimise the \emph{number} of queries'' -- halve the space;
  \item ``guarantee a win, but a wrong guess \emph{costs} you'' -- min-max, not midpoint (LC 375);
  \item ``rebuild a hidden tree using distance / ancestor queries'' -- adaptive group testing on
        the \chref{c24} oracle.
}
\recnotuse{
  \item the full array/matrix is visible up front -- \secref{c25-parent}/\secref{c25-degree};
  \item you just need to search a \emph{sorted} array once, not drive an interactive protocol
        -- ordinary binary search, \chref{c01};
  \item answers may be \emph{lies} (Ulam's game) -- that needs error-correcting splits, out of scope.
}
One-line test: \emph{can you only learn the answer by \textbf{asking}, and does each reply shrink
the possibilities? Choose the question that halves -- this section.}
\end{recognitionbox}

\subsection{Intuition}

You are on a number line $[1,n]$; the secret $x$ is somewhere on it and a referee only says
``higher'', ``lower'' or ``got it'' after each of your guesses. The clever move is always to guess
the \emph{middle} of the still-possible interval: whatever the referee says, you discard half the
remaining candidates. The interval of possibilities shrinks $n\to n/2\to n/4\to\cdots$, so about
$\log_2 n$ guesses finish it -- and you cannot do better, because each three-way answer can at best
divide the world into three and the count of distinguishable secrets grows only exponentially in the
number of questions. This is the whole game: \emph{a good query is one whose answers leave you with
equally small sub-cases}.

\begin{figure}[H]
\centering
\begin{tikzpicture}[diagram, x=11mm, y=8mm]
  % candidate interval rows, shrinking
  \node[annot, anchor=east] at (0,0) {\small $[1,\,10]$};
  \foreach \k in {1,...,10} { \node[cellnode, fill=primaryblue!10] at (\k,0) {\k}; }
  \node[annot, anchor=east] at (0,-1.1) {\small ask $5$};
  \foreach \k in {1,...,10} { \node[cellnode, fill=primaryblue!6] at (\k,-1.1) {\k}; }
  \fill[difforange!25] (0.62,-1.48) rectangle (5.38,-0.72);   % discard 1..5 (guess=5 correct-branch shown)
  \node[annot, anchor=east] at (0,-2.2) {\small ans.\! high};
  \foreach \k in {1,...,10} { \node[cellnode, fill=primaryblue!6] at (\k,-2.2) {\k}; }
  \fill[successgreen!28] (5.62,-2.58) rectangle (10.38,-1.82); % keep 6..10
  \node[annot, align=left, font=\scriptsize] at (5.5,-3.5)
    {guess the midpoint $5$; if told ``secret is higher'' the survivors $[6,10]$ are half\\
     of what remains -- each truthful answer throws away $\approx$ half the candidates};
\end{tikzpicture}
\caption{Adaptive halving on a number line. One well-chosen question (the midpoint) turns a size-$10$
candidate set into a size-$\le5$ one no matter which branch the oracle confirms. Repeating gives a
decision \emph{tree} of depth $\lceil\log_2 n\rceil$ -- exactly the shape the tree chapters draw, but
here its nodes are \emph{questions}, not data.}\label{fig:c25-halve}
\end{figure}

\subsection{Definitions and Notation}

\begin{definition}{Query (oracle) model}{c25-query}
A \term{query model} is a tuple $(\mathcal{X},\mathcal{Q},\mathit{ans})$: a hidden target
$x\in\mathcal{X}$ (the candidate space), a set of allowed \term{queries} $q\in\mathcal{Q}$, and an
answer function $\mathit{ans}(q,x)$ returning a feedback symbol. An \term{adaptive strategy} chooses
each next query as a function of the answers so far, and must halt with a unique $x$. Its
\term{cost} is the number (or weighted total) of queries in the worst case over $x$.
\end{definition}

\noindent\textbf{In plain words.} You can't see $x$; you may only poke questions and get terse
truthful replies; a strategy is good if it always names $x$ using few questions, and \emph{adaptive}
means your next question depends on the last answer.

\begin{proposition}{Halving and the decision-tree bound}{c25-halve}
If every query has at most $b$ possible answers, then $k$ queries can distinguish at most $b^{k}$
targets, so any correct adaptive strategy needs
$\displaystyle k\ \ge\ \lceil \log_{b} |\mathcal{X}| \rceil$ queries in the worst case. For a
three-way comparison oracle ($b=3$: higher/lower/equal) on $N$ ordered secrets this is
$\lceil\log_3 N\rceil$ \emph{in principle}, and binary search attains the (slightly stronger,
because answers must respect order) bound of $\lceil\log_2(N{+}1)\rceil$ guesses.
\end{proposition}
\begin{proofbox}
Each answer can at best multiply the number of distinguishable outcomes by $b$: after $k$ queries the
transcript of answers has $\le b^{k}$ distinct values, and a correct strategy must map different
targets to different transcripts, so $|\mathcal{X}|\le b^{k}$, i.e.\ $k\ge\log_b|\mathcal{X}|$. The
decision process is literally a $b$-ary \emph{tree} (nodes $=$ queries, edges $=$ answers, leaves $=$
conclusions) of depth $k$ and hence $\le b^{k}$ leaves. For an ordered search each comparison splits
the remaining interval into a left block, a point, and a right block; always splitting so the larger
side is as small as possible yields the recurrence $T(N)=1+T(\lfloor N/2\rfloor)$ whose solution is
$\lceil\log_2(N{+}1)\rceil$, attained by guessing the midpoint. \end{proofbox}

\begin{theorem}{Binary search is query-optimal}{c25-logn}
Against a truthful higher/lower/equal oracle, guessing the midpoint of the current candidate interval
minimises the worst-case number of queries; it finds any secret in $[1,n]$ within
$\lceil\log_2(n{+}1)\rceil$ guesses, matching \propref{c25-halve}'s lower bound up to the base. No
strategy can guarantee fewer in the worst case.
\end{theorem}
\begin{insightbox}
Optimality is \emph{about the split}, not the midpoint per se: you want the answer that leaves the
\emph{largest} remaining branch as small as possible. For a pure \emph{count} of questions on a
uniformly ordered line, that is the midpoint -- binary search. \emph{When answers cost something}
(LC 375: you pay the number you guessed) the split is no longer the middle -- you minimise the worst
case \emph{sum}, giving $k+\max(\text{left},\text{right})$ rather than $\max$ over sizes. Same
decision tree, different objective.
\end{insightbox}

\subsection{The algorithm: optimal play against a three-way oracle}

\begin{lstlisting}[caption={Guess Number Higher or Lower -- overflow-safe binary search on a comparison oracle}, label={lst:c25-guess}]
#include <bits/stdc++.h>
using namespace std;

// Forward-declared oracle (provided by the judge):
//   returns -1 if num is HIGHER than the secret ("too high"),
//            1 if num is LOWER than the secret ("too low"),
//            0 if num equals the secret.
int guess(int num);

class Solution {
public:
    int guessNumber(int n) {
        int lo = 1, hi = n;
        while (lo <= hi) {
            int mid = lo + (hi - lo) / 2;      // NOT (lo+hi)/2 -- lo+hi overflows int (n up to 2^31-1)
            int g = guess(mid);
            if (g == 0)  return mid;           // hit
            if (g == -1) hi = mid - 1;         // mid too high -> secret is in [lo, mid-1]
            else         lo = mid + 1;         // mid too low  -> secret is in [mid+1, hi]
        }
        return -1;                             // unreachable for a valid secret in [1,n]
    }
};
\end{lstlisting}
\begin{complexitybox}
Time $\bigO{\log n}$ guesses -- the interval halves every iteration (\thmref{c25-logn}); at
$n=2^{31}{-}1$ that is $\le 31$ oracle calls. Memory $\bigO{1}$. The \emph{overflow guard}
\code{lo + (hi-lo)/2} is the real content of an ``easy'' problem: with $n$ up to $2^{31}{-}1$, the
naive \code{(lo+hi)/2} exceeds \code{INT\_MAX} and mis-splits.
\end{complexitybox}

\subsection{Worked Example}

\begin{example}{Finding $x=2$ in $[1,10]$ by adaptive halving}{c25-query-run}
Secret $x=2$. Start $[\mathit{lo},\mathit{hi}]=[1,10]$.
Guess $\mathit{mid}=1{+}(10{-}1)/2=5$; $\mathit{guess}(5)=-1$ (too high) $\Rightarrow$
$\mathit{hi}=4$, interval $[1,4]$. Guess $\mathit{mid}=2$; $\mathit{guess}(2)=-1$ (too high)
$\Rightarrow$ $\mathit{hi}=1$, interval $[1,1]$. Guess $\mathit{mid}=1$; $\mathit{guess}(1)=1$
(too low) $\Rightarrow$ $\mathit{lo}=2$, and the loop's $\mathit{mid}=2$ gives $0$ after at most
$\lceil\log_2 11\rceil=4$ guesses. Every step threw away about half the remaining candidates -- the
\figref{c25-halve} shrink in action, finishing within the \thmref{c25-logn} bound.
\end{example}
\begin{tracebox}[Why the midpoint is the right question]
From a size-$m$ interval a midpoint guess leaves $\lfloor m/2\rfloor$ candidates on the ``too high''
branch and $\lceil m/2\rceil-1$ on the ``too low'' branch, so the larger survivor is $\le\lceil m/2\rceil$.
Guessing \emph{any other} $k$ makes one branch $k-1$ or $m-k$ candidates -- both can be $>m/2$ --
so no single guess has a smaller worst case. That is the \propref{c25-halve} lower bound met exactly.
\end{tracebox}

\subsection{Variations and Twists}

\begin{itemize}[nosep]
  \item \textbf{Cost-aware split (LC 375)}: you \emph{pay} the guessed number and want to minimise
        worst-case total cash; optimal first guess is not the middle -- interval DP over
        $[i,j]$. \hfill\lv{3}--\lv{4}
  \item \textbf{Reconstruct a hidden \emph{tree} by distance queries}: query $\dist(u,v)$ against an
        oracle (answerable via \chref{c24}); adaptively locate a node or edge by which split shrinks
        the candidate-tree set -- group testing on a tree. \hfill\lv{4}--\lv{5}
  \item \textbf{Ancestor/yes/no queries}: ``is $u$ an ancestor of $v$?'' answered in $\bigO{1}$ by
        $\tin/\tout$ (\thmref{c24-anc}) is a binary query; $n$ of them rebuild parent pointers by
        bisecting each node's ancestor chain. \hfill\lv{4}
  \item \textbf{Non-adaptive (batch) queries}: if all questions must be fixed upfront (no adaptation),
        you need combinatorial group designs -- more queries than the adaptive $\log$. \hfill\lv{5}
  \item \textbf{Lies allowed}: Ulam's game -- each answer might be wrong $e$ times; the candidate
        count is spread over $\sum_{i\le e}\binom{k}{i}$ transcripts, inflating queries. \hfill\lv{5}
\end{itemize}

\subsection{Common Pitfalls}\label{sec:c25-query-pitfalls}

\begin{warningbox}[Trap 1 -- \code{(lo+hi)/2} overflowing]
With $n$ up to $2^{31}{-}1$, \code{lo+hi} exceeds \code{INT\_MAX} and wraps negative, corrupting the
midpoint (wrong answer or an infinite loop). Always write \code{lo + (hi-lo)/2} -- mathematically
equal, overflow-free.
\end{warningbox}
\begin{lstlisting}[caption={\textcolor{warningred}{WRONG}: overflows when \code{lo+hi} exceeds \code{INT\_MAX}}, language=C++, basicstyle=\ttfamily\small]
int mid = (lo + hi) / 2;         // lo,hi up to 2^31-1  ->  signed overflow (UB)
\end{lstlisting}
\begin{lstlisting}[caption={\textcolor{successgreen}{RIGHT}: overflow-safe midpoint}, language=C++, basicstyle=\ttfamily\small]
int mid = lo + (hi - lo) / 2;    // hi-lo can't overflow; result identical
\end{lstlisting}

\begin{warningbox}[Trap 2 -- flipping the oracle's sign]
\code{guess(mid)}$=-1$ means \emph{your guess is too high}, so the secret is \emph{smaller}: shrink
$\mathit{hi}$, not $\mathit{lo}$. Swapping the two branches makes the interval grow and the loop
never ends (or returns wrong). Read the oracle's contract literally.
\end{warningbox}

\begin{warningbox}[Trap 3 -- non-shrinking bounds $\Rightarrow$ infinite loop]
Updating to $\mathit{hi}=\mathit{mid}$ (instead of $\mathit{mid}-1$) while $\mathit{lo}$ stays can keep
$\mathit{mid}$ constant forever. Because you already know $\mathit{mid}$ is \emph{not} the answer on
that branch, exclude it: $\mathit{hi}=\mathit{mid}-1$ / $\mathit{lo}=\mathit{mid}+1$.
\end{warningbox}

\subsection{Problem Lab}

LC 374 is the section's cleanest specimen: a hidden integer, a three-way oracle, and the provably
optimal halving strategy -- with one real trap (overflow) hiding behind an ``easy'' label.

% ---------------------------------------------------------------------------
\solvedproblem{c25-guess}{Guess Number Higher or Lower}{LeetCode 374}
  {https://leetcode.com/problems/guess-number-higher-or-lower/}{LC Easy L2}
  {binary search, interactive/oracle model, three-way comparison, overflow-safe midpoint}

\begin{problemstatement}[Guess Number Higher or Lower]
We pick a number $x$ in $[1, n]$. You guess a number and a pre-defined function \code{guess(num)}
returns \textbf{$-1$} if your guess is \emph{higher} than $x$ (``too high''), \textbf{$1$} if your guess
is \emph{lower} than $x$ (``too low''), and \textbf{$0$} if correct. Return the number $x$.
$1\le n\le 2^{31}-1$.
\end{problemstatement}

\begin{attemptgate}[Attempt it yourself first -- \lv{2}, give yourself 12 minutes]
Fair first attempt: obviously binary search. But $n$ is enormous -- make sure your loop terminates
and your midpoint never overflows.
\peek{What is the largest value \code{lo+hi} can take when $n=2^{31}-1$, and what does \code{int}
do with it?}
\nudge{Keep a $[\mathit{lo},\mathit{hi}]$ window. Compute \code{mid} as \code{lo+(hi-lo)/2}. On
\code{guess(mid)}$=-1$ move \code{hi} below \code{mid}; on $1$ move \code{lo} above \code{mid}; on $0$
return.}
\fullhint{Standard three-way binary search with an overflow-safe midpoint and correct branch
direction (\code{$-1$ $\Rightarrow$ secret smaller}). Loop while \code{lo<=hi}. Full build and proof
below.}
\end{attemptgate}

\solstep{Understanding the Problem}
A monotone predicate: define $P(k)\equiv(\mathit{guess}(k)\le 0)$, i.e.\ $k\ge x$. Truth grows with
$k$, so $x$ is the smallest $k$ with $P(k)$ true -- exactly binary search on a sorted yes/no boundary.
The oracle gives three-way feedback ($<,=,>$), so you can finish the moment $\mathit{guess}(\mathit{mid})=0$.
The only genuine difficulty is numeric: $n$ reaches \code{INT\_MAX}, so careless midpoint arithmetic
overflows.

\solstep{Key Observations}
\begin{observation}
Each answer collapses the candidate interval to at most half if you guess the midpoint
(\propref{c25-halve}), so $\le\lceil\log_2(n{+}1)\rceil$ guesses (about $31$ for the largest $n$) always
suffice -- and no strategy guarantees fewer (\thmref{c25-logn}).
\end{observation}
\begin{observation}
\code{guess(mid)}$=-1$ means the secret is \emph{strictly below} \code{mid}; $=1$ means \emph{strictly
above}. Because \code{mid} is known wrong on that branch, exclude it (\code{mid$\pm$1}) -- this both
prevents infinite loops and keeps the count at $\log$.
\end{observation}
\begin{observation}
\code{lo+hi} can overflow \code{int}; \code{lo+(hi-lo)/2} cannot, since $0\le\mathit{hi}-\mathit{lo}
\le n-1$ fits. This one line is the difference between AC and a Wrong Answer on large $n$.
\end{observation}

\solstep{Think It Through With Me}
\begin{thinkbox}
\deadend{Linear scan $1,2,\dots$ until \code{guess}$=0$.}{$O(n)$ oracle calls -- at $n=2^{31}$ it is
hopeless (TLE) even though each call is cheap. Queries are the currency here.}
\deadend{Guess \code{(lo+hi)/2}.}{On the biggest inputs \code{lo+hi} overflows to a negative \code{int},
\code{mid} lands outside $[\mathit{lo},\mathit{hi}]$, and you loop or return wrong. The idea is right;
the \emph{arithmetic} betrays it.}
\aha{Binary search with \code{mid=lo+(hi-lo)/2}: halve the window each answer, move to
\code{mid$\pm$1} (mid is known-wrong on that branch), stop when \code{guess=0}. Optimal and overflow-free.}
\end{thinkbox}

\solstep{Solution}
Maintain a closed candidate interval $[\mathit{lo},\mathit{hi}]=[1,n]$. While non-empty, compute the
overflow-safe midpoint, ask \code{guess}: on $0$ return it; on $-1$ set $\mathit{hi}=\mathit{mid}-1$
(secret lower); on $1$ set $\mathit{lo}=\mathit{mid}+1$ (secret higher). Monotonicity guarantees the
secret stays inside the window; halving guarantees termination in $\log$ steps.

\solstep{Proof of Correctness}
\begin{proofbox}
\emph{Invariant}: the secret $x\in[\mathit{lo},\mathit{hi}]$ at every loop start (true initially). If
$\mathit{guess}(\mathit{mid})=-1$ then $x<\mathit{mid}$, so setting $\mathit{hi}=\mathit{mid}-1$ keeps
$x$ in the window; the $1$ case is symmetric; the $0$ case returns $\mathit{mid}=x$. Thus the invariant
is preserved and the returned value equals $x$. \emph{Termination}: each iteration shrinks
$\mathit{hi}-\mathit{lo}$ to at most $\lceil(\mathit{hi}-\mathit{lo})/2\rceil-1<\mathit{hi}-\mathit{lo}$,
so after $\le\lceil\log_2(n{+}1)\rceil$ steps the window empties (and must have hit $0$ first, since
$x$ is always inside it). \emph{Optimality}: by \propref{c25-halve}, any strategy needs
$\Omega(\log n)$ three-way answers to distinguish $n$ ordered secrets, and \thmref{c25-logn} shows
this attains the bound, so it is query-optimal. \end{proofbox}

\solstep{Implementation}
See \coderef{c25-guess} above -- the listing in this section \emph{is} the accepted solution:
three-way binary search with \code{lo+(hi-lo)/2} and correct branch directions.

\solstep{Trace}
\begin{tracebox}[Official $n=10, x=6$]
$[\mathit{lo},\mathit{hi}]=[1,10]$: $\mathit{mid}=5$, $\mathit{guess}=1$ (low) $\Rightarrow\mathit{lo}=6$.
$[6,10]$: $\mathit{mid}=8$, $\mathit{guess}=-1$ (high) $\Rightarrow\mathit{hi}=7$.
$[6,7]$: $\mathit{mid}=6$, $\mathit{guess}=0$ $\Rightarrow$ return $6$. \checkmark\ Three guesses for
$x=6$ -- within $\lceil\log_2 11\rceil=4$.
\end{tracebox}

\solstep{Complexity}
\begin{complexitybox}
Time $\bigO{\log n}$ oracle calls ($\le 31$ at the maximum $n$); each call is treated as $O(1)$.
Space $\bigO{1}$. Optimal by \thmref{c25-logn}. Verdict: \ac.
\end{complexitybox}

\solstep{Pitfalls and Variants}
\begin{warningbox}
(i) \code{(lo+hi)/2} overflow (Trap~1) is the single most common WA/TLE here. (ii) Sign flip of the
oracle (Trap~2). (iii) Using $\mathit{hi}=\mathit{mid}$ / $\mathit{lo}=\mathit{mid}$ instead of
$\pm1$ (Trap~3) -- infinite loop on the boundary. \emph{Variant:} LC 374's ``guess the word'' (1095)
and ``find peak element'' (162) are the same adaptive halving on a different oracle -- binary search
without an explicit $=$ answer.
\end{warningbox}

\begin{insightbox}[Key takeaways]
\begin{enumerate}[nosep, leftmargin=1.6em]
  \item \textbf{Choose questions that halve.} With an oracle, your power is picking the query whose
        answers split the candidate space evenly (\propref{c25-halve}).
  \item \textbf{Counting bound $=$ decision tree.} $\log$ questions is a tree of depth $\log$ -- the
        same object as every other chapter, but nodes are \emph{queries}.
  \item \textbf{Easy label, real trap.} ``Easy'' hid an integer-overflow bug; correctness of the
        \emph{arithmetic} matters as much as the algorithm.
\end{enumerate}
\end{insightbox}

\begin{hardeningbox}
\begin{itemize}[nosep, leftmargin=1.4em]
  \item \textbf{Prove optimality}: show no decision tree with fewer than $\lceil\log_2(n{+}1)\rceil$
        leaves covers $n$ ordered secrets -- a counting argument (\propref{c25-halve}).
  \item \textbf{Cost-aware version}: replace the midpoint with the min-max choice
        $\arg\min_k\big(k+\max(T[\mathit{lo},k{-}1],T[k{+}1,\mathit{hi}])\big)$ -- LC 375.
  \item \textbf{Tree reconstruction by distance}: to find a hidden node $x$, repeatedly query
        $\dist(\mathit{centroid},x)$ (answerable in $O(1)$ after \chref{c24}) to discard whole
        subtrees -- centroid-halving instead of line-halving (\chref{c09} if centroid known).
\end{itemize}
\end{hardeningbox}

\subsection{Practice Problems}

\begin{stuckbox}
\begin{itemize}[nosep, leftmargin=1.4em]
  \item Did you write \code{lo+(hi-lo)/2}, not \code{(lo+hi)/2}? (Trap~1.)
  \item Did you map \code{guess}$=-1$ to ``search \emph{lower}'' (shrink \code{hi})? (Trap~2.)
  \item Did you exclude the known-wrong \code{mid} with $\pm1$? (Trap~3.)
  \item Can you state why fewer than $\log_2 n$ guesses cannot always work? (\propref{c25-halve}.)
\end{itemize}
\end{stuckbox}

\begin{practiceset}
  \practice{Guess Number Higher or Lower II}{LeetCode 375}
    {https://leetcode.com/problems/guess-number-higher-or-lower-ii/}{LC Medium L3}
    {Same three-way oracle as \probref{c25-guess}, but a wrong guess of $k$ \emph{costs you \$k} and
    you must return the minimum cash needed to \textbf{guarantee} a win regardless of the secret. The
    halving midpoint is no longer optimal -- you now minimise the worst-case \emph{total cost}, a
    different objective on the same decision tree. Model it with interval DP: let $T[i][j]$ be the
    guarantee for secrets $i..j$; then $T[i][j]=\min_{k\in[i,j]}\big(k+\max(T[i][k-1],\,T[k+1][j])\big)$
    with $T[i][j]=0$ when $i\ge j$ (the adversary picks the pricier branch, so you prepare for the
    $\max$). This is the ``cost-aware split'' variation from \secref{c25-query}; it also generalises
    the optimal-BST / matrix-chain family -- an interval-DP technique, not a tree-representation one.}
\end{practiceset}

\begin{recallbox}[10-second refresher -- 6.4 Recover by Queries (Interactive)]
You only see the answer by \emph{asking}: a query model is (space, allowed questions, feedback). A
truthful answer is worth a branch, so $k$ $b$-way answers distinguish $b^{k}$ targets -- you cannot
beat $\lceil\log_b\rceil$ (\propref{c25-halve}). Guess the midpoint: three-way binary search finds a
secret in $[1,n]$ in $\lceil\log_2(n{+}1)\rceil$ guesses, provably optimal (\thmref{c25-logn}), with
\code{lo+(hi-lo)/2} to dodge \code{INT\_MAX} overflow and $\mathit{mid}\pm1$ to keep shrinking. When
wrong guesses \emph{cost}, minimise $k+\max(\text{left},\text{right})$ over the interval (LC 375); when
the hidden object is a tree, halve with distance/ancestor queries off the \chref{c24} oracle.
\selftest{
  \item State the query model triple and give the oracle for LC 374.
  \item Prove $\lceil\log_2(n{+}1)\rceil$ is a lower bound on worst-case guesses for a three-way oracle.
  \item Why does \code{(lo+hi)/2} fail at $n=2^{31}-1$, and what replaces it?
  \item In LC 375, why do we take $\max$ (not sum/average) of the two sub-intervals plus the cost $k$?
  \item Give one distance/ancestor query that would help reconstruct a hidden tree, and why it halves.
}
\end{recallbox}

% !TEX root = ../../tree_notes.tex
% =====================================================================
%  Chapter 6 review -- the 15-section revision contract (chapter_end.txt).
%  Revision-only: tables, one-liners, pointers back via \thmref/\probref/
%  \secref/\coderef/\propref/\defref/\exref/\figref. No derivations.
%  Summit: LC Hard 1719 Number of Ways to Reconstruct a Tree
%  (given the ancestor-pairs, hunt a parent by degree and certify transitivity).
% =====================================================================

\chaptersection{Chapter 6 Summary}

Four sections, one inverse question: \emph{given raw data a tree left behind, rebuild the tree and
prove the rebuild is the tree}. The parent/edge view (6.1) turns pairs into adjacency and reads the
root off an in-degree histogram; the degree/code view (6.2) decides whether a degree list can be a
tree at all and counts the ones that are; the identification view (6.3) uses a union-find pre-look
to say \emph{tree, forest, or cycle} in one pass; the query view (6.4) rebuilds a hidden object one
adaptive question at a time. Every entry below is a pointer, not a re-derivation.

\begin{center}
\renewcommand{\arraystretch}{1.25}
\begin{tabular}{@{}p{0.25\linewidth} p{0.47\linewidth} p{0.20\linewidth}@{}}
\toprule
\textbf{Result} & \textbf{Statement (with conditions)} & \textbf{Where} \\
\midrule
Tree certificate & For a graph on $n$ nodes, any \textbf{two} of \{$n{-}1$ edges, connected,
  acyclic\} force the third; check two, never all three. & \thmref{c25-cert} \\
Degree sum & In every tree $\sum_v\degop(v)=2(n-1)$; the count leg of the certificate in disguise.
  & \thmref{c25-degsum} \\
Read the root & A rooted structure has exactly one node of in-degree $0$ (the root) and every other
  node of in-degree exactly $1$; two parents or two roots means \emph{not a tree}. & \propref{c25-root} \\
Degree-sequence realisability & Positive $d_1,\dots,d_n$ is tree-realisable \textbf{iff}
  $\sum d_i=2(n-1)$; the leaf-peel builds one. & \thmref{c25-degreal} \\
Pr\"ufer / count & Node $v$ appears $\degop(v)-1$ times in the length-$(n{-}2)$ code; the number of
  labelled trees with degrees $d_i$ is $(n-2)!\,/\,\prod_i(\degop(v_i)-1)!$. & \propref{c25-prufdeg} \\
Operational tree test & $n$ nodes form a tree \textbf{iff} there are $n-1$ edges \emph{and} union-find
  rejects none; $k$ rejections means $k$ surplus (cyclomatic) edges. & \corref{c25-istree} \\
Query-optimal search & Against a truthful higher/lower/equal oracle the midpoint split attains
  $\lceil\log_2(n{+}1)\rceil$ guesses, and no strategy is better in the worst case. & \thmref{c25-logn} \\
\bottomrule
\end{tabular}
\end{center}

\medskip
{\small\textbf{Algorithms from this chapter -- rebuild, then certify.}}
\begin{center}
\small
\renewcommand{\arraystretch}{1.2}
\setlength{\tabcolsep}{3pt}
\begin{tabular}{@{}p{0.19\linewidth} l l p{0.29\linewidth} p{0.155\linewidth}@{}}
\toprule
\textbf{Algorithm} & \textbf{Time} & \textbf{Memory} & \textbf{When to use (deciding phrase)} & \textbf{Code} \\
\midrule
Rebuild adj, find root, certify
  & $O(n)$ & $O(n)$
  & ``edge list / parent array given, is it a valid rooted tree?''
  & \coderef{c25-build} \\
Degree-realise by leaf-peel
  & $O(n\log n)$ & $O(n)$
  & ``build \emph{a} tree with these degrees / is this degree list a tree?''
  & \coderef{c25-degbld} \\
Union-find pre-look
  & $O((n{+}m)\,\alpha)$ & $O(n)$
  & ``is it a tree? / how many components? / delete the redundant edge''
  & \coderef{c25-uf} \\
Optimal guess (binary search)
  & $O(\log n)$ & $O(1)$
  & ``find a hidden value through a higher/lower oracle, fewest questions''
  & \coderef{c25-guess} \\
Summit: degree-greedy parent hunt
  & $O(n^3)$ & $O(n^2)$
  & ``given all ancestor pairs, how many rooted trees match?''
  & \coderef{c25-summit} \\
\bottomrule
\end{tabular}
\end{center}

\chaptersection{Master Decision Flowchart}

Post-study ladder: walk top to bottom, stop at the first branch that matches. It decides
\emph{which pile of data you were handed} and hence which rebuild-and-certify tool to run -- a
follower never trusts a count that was never a certificate.

\begin{lstlisting}[language={}, basicstyle=\ttfamily\small, numbers=none, frame=single]
  You are handed DATA about a tree. What shape is the data?
  |
  |-- an EDGE LIST or PARENT ARRAY (the pairs are explicit)
  |     +-- undirected, exactly n-1 edges -> count leg + one union-find pass   [6.3]
  |     +-- rooted / directed -> read in-degree: one 0 (root), rest exactly 1  [6.1, LC 685]
  |     +-- n edges, one on a cycle -> delete the LAST rejected edge           [6.3, LC 684]
  |     +-- raw equation / relationship pairs -> rebuild adj, DFS the ratio    [6.1, LC 399]
  |
  |-- a DEGREE SEQUENCE or CODE (only statistics survive)
  |     +-- all d_i >= 1 and sum(d) == 2(n-1)?  else NOT realisable           [6.2]
  |     +-- yes -> leaf-peel (Pruefer decode) builds one tree                  [6.2]
  |     +-- "how many trees?" -> (n-2)! / prod (d_i - 1)!  (Cayley = all d free)[6.2]
  |     +-- grid / region uniformity -> recurse into 4 quadrants (quad tree)   [6.2, LC 427]
  |
  |-- "IS IT a tree / how many components?" (cycle & connectivity ask)
  |     +-- union-find: a rejected edge = a cycle; components = final set count [6.3]
  |     +-- directed acyclicity (order exists) -> topological sort, NOT UF     [6.3, LC 210]
  |
  |-- an ORACLE (you only get to ASK)
  |     +-- higher/lower/equal -> binary search, ceil(log2(n+1)) is optimal     [6.4, LC 374]
  |     +-- a wrong guess COSTS you -> min-max k + max(left,right), interval DP  [6.4, LC 375]
  |
  +-- recursion / matrix sized by n -> a chain input is O(n) deep: explicit stack,
      or a vector<vector<bool>> sized n^2 when an adj list would do. Full DSU &
      MST -> [c17]; Pruefer/Cayley proofs -> [c32]; distance/LCA oracle -> [c24].
\end{lstlisting}

\chaptersection{Common Mistakes Checklist}

\begin{itemize}[nosep, leftmargin=1.4em]
  \item \textbf{Trusting $n-1$ edges alone.} A disconnected forest can carry $n-1$ edges. The
        certificate needs \emph{two} of the three legs -- pair the count with connectivity or with
        a cycle-free union-find pass, never the count by itself (\thmref{c25-cert}, Trap~1).
  \item \textbf{Reading in-degree on an undirected list.} Feeding both directions makes the root
        look like it has a parent and every leaf look like a non-root -- a rooted instance (LC 685)
        must scan \emph{directed} edges (\propref{c25-root}, Trap~2).
  \item \textbf{Banning the wrong edge in the two-parent repair.} When a node has two parents,
        \emph{remove the later edge first}, then run the cycle scan on the rest; returning the
        second parent without testing the cycle mis-picks the answer (\probref{c25-rc2}).
  \item \textbf{Confusing the degree \emph{sum} with a \emph{realisation}.} $\sum d=2(n-1)$ with
        all $d_i\ge1$ is necessary \emph{and} sufficient -- but a $0$ or a negative sneaking in
        breaks the peel; screen for it first (\thmref{c25-degreal}, Trap~1).
  \item \textbf{Inverting the Pr\"ufer count.} Node $v$ appears $\degop(v)-1$ times \emph{in the
        code}, and the count divides by $\prod(\degop(v)-1)!$ -- not $\degop(v)!$
        (\propref{c25-prufdeg}, Trap~2).
  \item \textbf{Running union-find on a directed acyclicity ask.} ``is a topological order
        possible?'' (LC 210) is \emph{not} the undirected cycle test; union-find ignores edge
        direction and gives the wrong verdict (\secref{c25-identify}, Trap~3).
  \item \textbf{Off-by-one and overflow in the oracle.} The interval is $[\text{lo},\text{hi}]$
        closed -- move to \code{mid-1}/\code{mid+1}; and \code{(lo+hi)/2} overflows for
        $n$ up to $2^{31}{-}1$, so use \code{lo+(hi-lo)/2} (\probref{c25-guess}, Trap~1).
\end{itemize}

\chaptersection{WA Clinic}

Verdicts are symptoms; the causes are the short list this chapter drilled. Diagnose top-down --
every leaf names the training problem.

\begin{figure}[H]
\centering
\begin{tikzpicture}[diagram,
    v/.style={boxnode, fill=primaryblue!12, draw=primaryblue, text width=22mm, align=center},
    b/.style={boxnode, fill=difforange!12, draw=difforange, text width=12mm},
    l/.style={boxnode, fill=successgreen!10, draw=successgreen, text width=58mm,
              font=\scriptsize, align=left}]
  \node[v] (verdict) {\textbf{Verdict on the judge}};
  \node[b, below left=8mm and 6mm of verdict] (wa) {\textbf{WA}};
  \node[b, below right=1mm of verdict] (tle) {\textbf{TLE}};
  \node[b, below left=8mm and 1mm of verdict, xshift=30mm] (re) {\textbf{RE}};
  \node[b, below right=8mm and 6mm of verdict] (mle) {\textbf{MLE}};
  \node[l, below right=1mm and 2mm of wa.south] (wa1)
    {count-only ``tree'' flags a disconnected forest (\corref{c25-istree}); in-degree scan on an
     undirected list invents two parents (\probref{c25-rc2}); wrong banned edge (\probref{c25-rc2});
     Pr\"ufer count divides by $\degop!$ not $(\degop-1)!$ (\propref{c25-prufdeg}); oracle interval
     off by one (\probref{c25-guess})};
  \node[l, below right=1mm and 2mm of tle.south] (tle1)
    {DFS from scratch per candidate edge instead of one union-find pass (\probref{c25-rc1});
     $O(n^3)$ transitivity scan is fine at $n\le500$ but re-DFSing the graph per pair is not
     (\probref{c25-summit})};
  \node[l, below right=1mm and 2mm of re.south] (re1)
    {recursive quad-tree / tree-from-degrees build on a degenerate $n$-deep input -- go explicit
     or raise stack (\secref{c25-degree}, Trap~3); \code{(lo+hi)/2} overflow in the oracle
     (\probref{c25-guess})};
  \node[l, below right=1mm and 2mm of mle.south] (mle1)
    {an $n\times n$ \code{bool} matrix where an adjacency list suffices -- the Summit's
     $\mathit{adj}[\cdot][\cdot]$ is $O(n^2)$ on purpose at $n\le500$, but copying it per node is
     not (\probref{c25-summit})};
  \draw[flow] (verdict) -- (wa); \draw[flow] (verdict) -- (tle);
  \draw[flow] (verdict) -- (re); \draw[flow] (verdict) -- (mle);
  \draw[flow] (wa) -- (wa1); \draw[flow] (tle) -- (tle1);
  \draw[flow] (re) -- (re1); \draw[flow] (mle) -- (mle1);
\end{tikzpicture}
\caption*{Chapter~6 verdict tree -- only the failures this chapter's problems can produce.}
\end{figure}

\medskip
\noindent Backup table (same content, terse):
\begin{center}
\small
\renewcommand{\arraystretch}{1.2}
\begin{tabular}{@{}llp{0.42\linewidth}l@{}}
\toprule
\textbf{Verdict} & \textbf{Usual cause} & \textbf{One-line fix} & \textbf{Trained in} \\
\midrule
WA  & count-only ``is a tree''   & pair $n{-}1$ edges with a UF pass        & \corref{c25-istree} \\
WA  & two-parent false alarm     & scan \emph{directed} edges only          & \probref{c25-rc2} \\
WA  & Pr\"ufer count $\degop!$    & divide by $(\degop(v)-1)!$               & \propref{c25-prufdeg} \\
WA  & oracle off-by-one          & closed $[\text{lo},\text{hi}]$, move $\pm1$ & \probref{c25-guess} \\
TLE & re-DFS per candidate edge  & one union-find sweep                      & \probref{c25-rc1} \\
RE  & dense $\mathit{adj}$ matrix   & size once at $n\le500$; explicit stack & \probref{c25-summit} \\
\bottomrule
\end{tabular}
\end{center}

\chaptersection{Quick Reference Card}

\begin{quickref}[Chapter 6 at a glance]
\textbf{One question.} The rest of the book \emph{walks} a tree someone drew for you; this chapter
answers the inverse -- \emph{what data pins a tree down, and how do I rebuild and certify it} from
that data alone? Rebuild the object, then prove it is a tree with the certificate (\thmref{c25-cert}).

\medskip
\textbf{Two invariants, always.} $m=n-1$ edges (equivalently $\sum_v\degop(v)=2(n-1)$,
\thmref{c25-degsum}) and \emph{connected} (or \emph{acyclic}) -- any two of the three certificate
legs force the third (\thmref{c25-cert}). A degree list is realisable iff every $d_i\ge1$ and the
sum is exactly $2(n-1)$ (\thmref{c25-degreal}).

\medskip
\textbf{Four shadows of the same tree $T$.} Edge list / parent array $\to$ rebuild and read the
root (\secref{c25-parent}, \exref{c25-parent-run}). Degree sequence $\to$ peel and count (\secref{c25-degree},
\exref{c25-deg-run}). Component count $\to$ union-find rejections (\secref{c25-identify},
\exref{c25-id-run}). Oracle answers $\to$ adaptive halving (\secref{c25-query}, \exref{c25-query-run}).
All four describe the \emph{identical} $T$ of \figref{c25-running}.

\medskip
\textbf{Pr\"ufer one-liner.} Node $v$ shows up $\degop(v)-1$ times in the length-$(n{-}2)$ code, so
a degree sequence is realised by $(n-2)!\,/\,\prod_v(\degop(v)-1)!$ labelled trees
(\propref{c25-prufdeg}) -- Cayley's $n^{\,n-2}$ is that sum with degrees free (\chref{c32}).
\end{quickref}

\begin{lstlisting}[caption={Chapter 6 code kit: rebuild, realise, union-find, guess}, label={lst:c25-kit}]
#include <bits/stdc++.h>
using namespace std;

// (1) Edge list / parent array -> adjacency, read the root, certify.   [6.1]
vector<vector<int>> buildFromParent(const vector<int>& par, int& root, bool& ok) {
    int n = (int)par.size() - 1; vector<vector<int>> adj(n + 1); vector<int> indeg(n + 1);
    root = 0; int zeros = 0; ok = true;
    for (int v = 1; v <= n; ++v) { int p = par[v];
        if (p == 0) { root = v; ++zeros; }                 // candidate root(s)
        else if (p >= 1 && p <= n) { adj[p].push_back(v); ++indeg[v]; }
        else ok = false; }                                 // parent out of range
    if (zeros != 1) ok = false;                            // exactly one root
    for (int v = 1; v <= n; ++v) if (v != root && indeg[v] != 1) ok = false;
    return adj;                                            // + DFS from root for connectivity
}

// (2) Degree sequence -> a realising tree by the leaf-peel (Pruefer decode). [6.2]
bool realisable(const vector<int>& deg) {                  // the WHOLE test (Thm 6.2)
    long long sum = 0;
    for (int d : deg) if (d < 1) return false; else sum += d;
    return sum == 2LL * ((int)deg.size() - 1);
}

// (3) Union-find pre-look: a rejected edge closes a cycle; count components. [6.3]
struct DSU { vector<int> p;
    DSU(int n) : p(n) { iota(p.begin(), p.end(), 0); }
    int find(int x) { return p[x] == x ? x : p[x] = find(p[x]); }
    bool uni(int a, int b) { a = find(a); b = find(b);
        if (a == b) return false; p[b] = a; return true; } };   // false => cycle

// (4) Optimal guess against a higher/lower/equal oracle -- overflow-safe.    [6.4]
int guessNumber(int n) {
    int lo = 1, hi = n;
    while (lo <= hi) {
        int mid = lo + (hi - lo) / 2;                      // NOT (lo+hi)/2 -- overflow
        int g = guess(mid);                                // -1 too high, 1 too low, 0 hit
        if (g == 0) return mid;
        if (g == -1) hi = mid - 1; else lo = mid + 1;
    }
    return -1;
}
\end{lstlisting}

\chaptersection{Write-From-Memory Drills}

Five signature skeletons, each lifted from a listing already in the chapter. Cover the key
line, reproduce it from memory, then check the answer one line below.

\begin{lstlisting}[caption={Drill 1 -- rebuild a rooted tree from its parent array (from \coderef{c25-build})}]
for (int v = 1; v <= n; ++v) { int p = par[v];
    if (p == 0) { root = v; ++zeros; }
    else if (1 <= p && p <= n) { /* FILL: ______ */; ++indeg[v]; }
    else ok = false;
}
if (/* FILL: zeros != 1 */) ok = false;                    // exactly one root
\end{lstlisting}
\noindent{\scriptsize Key: append $v$ to \code{adj[p]} for the child list; validity is
\emph{one} zero in-degree and in-degree $1$ for every other node (\propref{c25-root}).}

\begin{lstlisting}[caption={Drill 2 -- the degree-sum gate (from \coderef{c25-degbld})}]
long long sum = 0;
for (int d : deg) { if (/* FILL: d < 1 */) return {}; sum += d; }
if (/* FILL: sum != 2LL*(n-1) */) return {};               // Thm 6.2 is the WHOLE test
\end{lstlisting}
\noindent{\scriptsize Key: positivity first, then the single equation $\sum d=2(n-1)$
(\thmref{c25-degreal}); if it passes, the leaf-peel never gets stuck.}

\begin{lstlisting}[caption={Drill 3 -- the Pr\"ufer multiset (from \coderef{c25-degbld})}]
for (int i = 0; i < n; ++i)
    for (int k = 1; k < /* FILL: deg[i] */; ++k) prufer.push_back(i);   // deg(v)-1 copies
\end{lstlisting}
\noindent{\scriptsize Key: node $v$ appears $\degop(v)-1$ times, so the code has length
$\sum(\degop(v)-1)=n-2$ (\propref{c25-prufdeg}).}

\begin{lstlisting}[caption={Drill 4 -- is-it-a-tree by one union-find pass (from \coderef{c25-uf})}]
bool isTree(int n, const vector<pair<int,int>>& edges) {
    if ((int)edges.size() != /* FILL: n - 1 */) return false;         // leg (1)
    DSU d(n);
    for (auto& e : edges) if (!d.uni(e.first, e.second)) return false; // a reject = a cycle
    return true;
}
\end{lstlisting}
\noindent{\scriptsize Key: $n-1$ edges \emph{and} no rejected edge $\Rightarrow$ connected and
acyclic at once (\corref{c25-istree}).}

\begin{lstlisting}[caption={Drill 5 -- optimal play against the oracle (from \coderef{c25-guess})}]
int lo = 1, hi = n;
while (lo <= hi) {
    int mid = /* FILL: lo + (hi - lo)/2 */;
    int g = guess(mid);
    if (g == 0) return mid;
    if (g == -1) /* FILL: hi = mid - 1 */; else /* FILL: lo = mid + 1 */;
}
\end{lstlisting}
\noindent{\scriptsize Key: closed interval, so shrink to \code{mid-1}/\code{mid+1}; the midpoint
guard avoids overflow at $n=2^{31}{-}1$ (\thmref{c25-logn}, Trap~1).}

\chaptersection{Complete Problem Checklist}

Every problem this chapter offers -- 4 fully solved in-section, 4 practice, one Summit, five
mixed. Tick \emph{only} after an unaided solve; a hinted attempt counts as half. Cite by name
and number -- URLs appear once, in the solved and practice entries.

\begin{center}
\footnotesize
\renewcommand{\arraystretch}{1.25}
\begin{tabular}{@{}c p{0.17\linewidth} l p{0.20\linewidth} p{0.24\linewidth} c@{}}
\toprule
\textbf{\#} & \textbf{Problem} & \textbf{Source} & \textbf{Key technique} & \textbf{What to look for} & \textbf{Done?} \\
\midrule
\multicolumn{6}{l}{\textbf{Section 6.1 -- parent array / edge list}} \\
W6.1.1 & Redundant Connection II & LC 685 & \lv4 in-degree + union-find repair & two-parent case, ban later edge, then cycle scan & $\square$ \\
P6.1.2 & Evaluate Division & LC 399 & \lv3 rebuild graph, ratio DFS & pour equation pairs into weighted adj, walk a path & $\square$ \\
\midrule
\multicolumn{6}{l}{\textbf{Section 6.2 -- degrees, sequences, codes}} \\
W6.2.1 & Construct Quad Tree & LC 427 & \lv3 recursive uniform partition & all-equal test then split into four quadrants & $\square$ \\
P6.2.2 & Course Schedule II & LC 210 & \lv3 topological order (directed) & prerequisite pairs $\to$ an ordering, or the empty list on a cycle & $\square$ \\
\midrule
\multicolumn{6}{l}{\textbf{Section 6.3 -- is-it-a-tree / union-find}} \\
W6.3.1 & Redundant Connection & LC 684 & \lv3 union-find first reject & $n$ edges, return the last edge that closed a cycle & $\square$ \\
P6.3.2 & Number of Provinces & LC 547 & \lv2 count components & matrix $\to$ union-find, answer is the final set count & $\square$ \\
\midrule
\multicolumn{6}{l}{\textbf{Section 6.4 -- queries / interactive}} \\
W6.4.1 & Guess Number Higher or Lower & LC 374 & \lv2 binary search on oracle & overflow-safe midpoint, closed-interval shrink & $\square$ \\
P6.4.2 & Guess Number Higher or Lower II & LC 375 & \lv4 min-max interval DP & wrong guess costs $k$: $k+\max(\text{lo},\text{hi})$ & $\square$ \\
\midrule
\multicolumn{6}{l}{\textbf{Summit (chapter end)}} \\
S6 & Number of Ways to Reconstruct a Tree & LC 1719 & \lv5 degree-greedy parent hunt + transitivity & equal-degree tie $\to 2$, missing containment $\to 0$ & $\square$ \\
\bottomrule
\end{tabular}
\end{center}

\chaptersection{Recommended Study Order}

Five focused days. Day 1 fixes the certificate and the rebuild; day 2 the degree/count machinery;
day 3 the union-find identification; day 4 the query model; day 5 the Summit. Solve in this order
even where it differs from section order -- each day's tool is already in hand the day before.

\begin{center}
\small
\renewcommand{\arraystretch}{1.25}
\begin{tabular}{@{}llp{0.55\linewidth}@{}}
\toprule
\textbf{Day} & \textbf{Focus} & \textbf{Tasks} \\
\midrule
1 & Read + do 6.1 & rebuild $T$ from parent array and edge list (\exref{c25-parent-run}); LC 685; LC 399; write the root-read cold \\
2 & Read + do 6.2 & realise $T$'s degree sequence by the peel (\exref{c25-deg-run}); LC 427; LC 210; count the $180$ realisations \\
3 & Read + do 6.3 & classify $T$ and $T$+bad edge (\exref{c25-id-run}); LC 684; LC 547 \\
4 & Read + do 6.4 & halve $[1,10]$ to find $x$ (\exref{c25-query-run}); LC 374; LC 375 \\
5 & Summit + revise & LC 1719 degree-greedy hunt; walk the decision flowchart on five questions \\
\bottomrule
\end{tabular}
\end{center}

\chaptersection{Difficulty Progression}

Tiers are competency gates -- each row names the \emph{skill} the tier unlocks, not just a
rating bucket.

\begin{center}
\small
\renewcommand{\arraystretch}{1.3}
\begin{tabular}{@{}lp{0.36\linewidth}p{0.34\linewidth}@{}}
\toprule
\textbf{Tier} & \textbf{Problems} & \textbf{What passing this tier proves} \\
\midrule
$\lv2$ \ $\bigstar\bigstar$
  & LC 374, 547
  & you can drive an oracle optimally and count components with one union-find pass \\
$\lv3$ \ $\bigstar\bigstar\bigstar$
  & LC 684, 427, 399, 210
  & rebuild from raw pairs, spot the redundant edge, and partition a grid recursively \\
$\lv4$ \ $\bigstar\bigstar\bigstar\bigstar$
  & LC 685, 375
  & split on in-degree \emph{and} cycle in a directed repair; shift the split when a guess costs \\
$\lv5$ \ $\bigstar\bigstar\bigstar\bigstar\bigstar$
  & LC 1719 (Summit)
  & reconstruct from partial (ancestor) data and decide existence and multiplicity \\
\bottomrule
\end{tabular}
\end{center}

\chaptersection{Reduction Gallery}

Construction is a disguise; each line collapses a family onto one rebuild-and-certify move. No new
theory -- just where the certificate is already hiding.

\begin{itemize}[nosep, leftmargin=1.4em]
  \item ``is this a tree?'' $\longrightarrow$ $n-1$ edges plus one union-find pass with no rejection
        (\corref{c25-istree}).
  \item ``delete one edge to make it a tree'' (LC 684) $\longrightarrow$ the last edge union-find
        rejects is on the unique cycle (\probref{c25-rc1}).
  \item ``one node has two parents / a cycle too'' (LC 685) $\longrightarrow$ in-degree scan splits
        the two repair cases before the cycle test (\probref{c25-rc2}).
  \item ``build \emph{a} tree with these degrees'' (LC-ish) $\longrightarrow$ check $\sum d=2(n-1)$,
        then peel leaves off a Pr\"ufer multiset (\thmref{c25-degreal}, \coderef{c25-degbld}).
  \item ``how many trees share this degree sequence?'' $\longrightarrow$ decode count
        $(n-2)!/\prod(\degop(v)-1)!$ (\propref{c25-prufdeg}).
  \item ``is this $n\times n$ grid uniform?'' (LC 427) $\longrightarrow$ a tree built from the
        data by four-way recursion, not from edges (\probref{c25-quad}).
  \item ``can I order these prerequisites?'' (LC 210) $\longrightarrow$ a \emph{directed} acyclicity
        test (topological sort), the cousin union-find cannot answer (\secref{c25-identify}).
  \item ``group accounts sharing an email'' (LC 721) $\longrightarrow$ union-find where the keys are
        strings -- the same merge mechanic as \coderef{c25-uf}.
  \item ``are these equality/inequality constraints satisfiable?'' (LC 990) $\longrightarrow$ union
        the equals, then any inequality whose two ends are already joined is a contradiction -- a
        cycle in disguise (\corref{c25-istree}).
  \item ``rebuild an N-ary tree inside a binary one'' (LC 431) $\longrightarrow$ first-child /
        next-sibling encoding, the structural mirror of \secref{c25-parent}'s parent-array rebuild.
  \item ``how many rooted trees match these ancestor pairs?'' (LC 1719) $\longrightarrow$ degree-based
        parent hunt plus transitivity containment (\probref{c25-summit}).
\end{itemize}

\chaptersection{Mixed Problem Set}

These refuse to name their pattern; decide \emph{which rebuild-and-certify move} you need before
opening the code. A hint tells you the shape of the answer, never the tool.

\begin{mixedset}
  \practice{Course Schedule}{LeetCode 207}
    {https://leetcode.com/problems/course-schedule/}{LC Medium L3}
    {``can you finish all courses'' is ``is the prerequisite digraph acyclic'' -- but note the
    direction matters, so this is the \emph{topological} twin of \secref{c25-identify}, not the
    undirected union-find test of \corref{c25-istree}: a Kahn/DFS in-degree peel decides it.}
  \practice{Regions Cut By Slashes}{LeetCode 959}
    {https://leetcode.com/problems/regions-cut-by-slashes/}{LC Hard L4}
    {scale the grid to a cell complex and union the four neighbours of every open cell; the answer
    is the \emph{component count} of \coderef{c25-uf} after you stop merging -- the exact
    ``components = final set count'' reading of \secref{c25-identify}.}
  \practice{Accounts Merge}{LeetCode 721}
    {https://leetcode.com/problems/accounts-merge/}{LC Medium L3}
    {emails that co-occur on one account belong to one person: union-find where the universe is
    email \emph{strings}, not indices. Group by the \term{find} representative -- the same merge
    mechanic as \coderef{c25-uf}, keyed differently.}
  \practice{Satisfiability of Equality Equations}{LeetCode 990}
    {https://leetcode.com/problems/satisfiability-of-equality-equations/}{LC Medium L3}
    {union every ``$a==b$'' pair, then every ``$a!=b$'' whose ends are \emph{already} joined is a
    contradiction -- an inequality edge closing a cycle over the equalities, so it is
    \corref{c25-istree}'s rejection test in a fresh costume.}
  \practice{Encode N-ary Tree to Binary Tree}{LeetCode 431}
    {https://leetcode.com/problems/encode-n-ary-tree-to-binary-tree/}{LC Hard L4}
    {a lossless encoding \emph{is} a reconstruction: map each node to \emph{first child} and
    \emph{next sibling} pointers so the N-ary tree survives as a binary one, then rebuild it exactly
    as \secref{c25-parent} rebuilds $T$ from a parent array -- the arrow reversed.}
  \practice{Remove Max Number of Edges to Keep Graph Fully Traversable}{LeetCode 1579}
    {https://leetcode.com/problems/remove-max-number-of-edges-to-keep-graph-fully-traversable/}{LC Hard L4}
    {Type-3 edges serve both players, so claim them last: union the type-1 owners, then
    the type-2, then spend \coderef{c25-uf}'s cycle test on the shared edges -- the
    maximal-removal question is \secref{c25-identify}'s "which edges are forced" run
    backwards.}
  \practice{Lexicographically Smallest Equivalent String}{LeetCode 1061}
    {https://leetcode.com/problems/lexicographically-smallest-equivalent-string/}{LC Medium L3}
    {The equalities partition 26 letters into classes; rewrite every character as its
    class's \emph{smallest} member -- \coderef{c25-uf}'s canonical-representative
    habit, one find per letter.}
  \practice{Most Stones Removed with Same Row or Column}{LeetCode 947}
    {https://leetcode.com/problems/most-stones-removed-with-same-row-or-column/}{LC Medium L3}
    {Every removal leaves the component connected, so what survives is one stone per
    component over shared row/column coordinates: union, then count the final
    representatives -- \secref{c25-identify}'s component ledger with stones as
    vertices.}
  \practice{Possible Bipartition}{LeetCode 886}
    {https://leetcode.com/problems/possible-bipartition/}{LC Medium L3}
    {Disliked pairs must land on opposite sides: BFS two-colouring answers directly, or
    union each person with all of a rival's rivals -- a failed merge is
    \corref{c25-istree}'s rejection test wearing an odd cycle.}
  \practice{Count the Number of Complete Components}{LeetCode 2685}
    {https://leetcode.com/problems/count-the-number-of-complete-components/}{LC Medium L3}
    {Group vertices with \coderef{c25-uf}, then tally vertices and edges per component:
    complete exactly when the edges hit $v(v-1)/2$. The certify-after-building step is
    \secref{c25-degree}'s degree audit counted on edges.}
\end{mixedset}


\chaptersection{The Summit -- LC 1719 Number of Ways to Reconstruct a Tree}

\begin{notebox}[Why this is the chapter's top rope]
Sections 6.1--6.4 each rebuild or verify a tree from \emph{one clean shadow}: a full edge list, a
degree sequence, a bag of union-find answers, an oracle. The Summit hands you the faintest shadow
imaginable -- only the set of \emph{ancestor--descendant pairs}, and not even the tree -- then asks
two questions at once: does \emph{any} rooted tree realise exactly these pairs, and if so are there
\emph{several}? The whole chapter converges here: you must \emph{reconstruct} a structure from
partial data (\secref{c25-parent}), reason about \emph{degrees} to guess who can be a parent
(\secref{c25-degree}), and \emph{certify} validity -- here by the transitivity of ancestry rather
than a union-find cycle count (\secref{c25-identify}). A node's candidate parent is the
\emph{lowest-degree node still connected to it}: ancestry is transitive, so every neighbour of a
node must also be a neighbour of its parent. When two nodes can occupy the same parent slot (equal
degree) the answer is ``more than one''; when a required containment fails, ``none.''
\end{notebox}

\solvedproblem{c25-summit}{Number of Ways to Reconstruct a Tree}{LeetCode 1719}
  {https://leetcode.com/problems/number-of-ways-to-reconstruct-a-tree/}{LC Hard L5}
  {degree-sorted parent hunt + ancestor-transitivity containment test; equal-degree tie $\to$ multiple}

\begin{problemstatement}[Number of Ways to Reconstruct a Tree]
You are given \code{pairs}, a list of $\text{pairs}[i]=[x_i,y_i]$ with $x_i<y_i$ and no duplicates,
describing an unknown rooted tree on the nodes that appear. A pair $[x,y]$ is present \textbf{iff}
one of $x,y$ is an \emph{ancestor} of the other (nodes on the root-to-node path, the node itself
excluded). The tree is \emph{not} binary -- a node may have any number of children. Return
\textbf{0} if no rooted tree matches, \textbf{1} if exactly one does, and \textbf{2} if more than
one does. Two trees differ when some node has a different parent. ($1\le|\text{pairs}|\le2\cdot10^{5}$,
$1\le x_i<y_i\le500$.)
\end{problemstatement}

\begin{attemptgate}[Attempt it yourself first -- \lv{5}, give yourself 50 minutes]
A fair first attempt: notice that a node connected to \emph{every} other node must be the root (its
ancestor set is everyone), and that if you fix a root, a node's parent is \emph{some} ancestor. But
enumerating trees is exponential -- the challenge is to decide existence and multiplicity
\emph{greedily} from the degrees alone.
\peek{sort the nodes by how many pairs they appear in. Who is the most likely parent of a node $x$
-- a node of larger or smaller degree?}
\nudge{a parent must be connected to all of $x$'s neighbours too (transitivity), so look for
$x$'s parent among nodes \emph{connected to} $x$ with degree at least $\degop(x)$. The lowest such
degree is the tightest fit -- the nearest ancestor.}
\fullhint{Build the adjacency and degrees. Process nodes in \emph{increasing} degree; for node $x$
scan the higher-degree nodes for the first one $y$ that is \emph{adjacent} to $x$ -- that is
$x$'s parent candidate. If $\degop(x)=\degop(y)$ the two could swap parent slots, so mark
\emph{multiple}. Verify $y$ is adjacent to \emph{every} neighbour of $x$ (containment); if any
fails, no tree exists ($0$). If more than one node ends up with no parent, more than one root
($0$). Else return $2$ if a tie was seen, otherwise $1$. Full build and proof below.}
\end{attemptgate}

\solstep{Understanding the Problem}
The input is not edges of the tree -- it is the \emph{ancestor relation} compressed into pairs. So
$[x,y]\in\text{pairs}$ does not mean $x$ is $y$'s parent, only that one is somewhere above the
other. Two structural facts drive everything. First, ancestry is \emph{transitive}: if $p$ is
$x$'s parent, then every node related to $x$ (its ancestors, and, going down, its descendants that
reach through $x$) must also be related to $p$ -- in the pair graph, $p$'s neighbour set must
\emph{contain} $x$'s neighbour set. Second, a parent has \emph{at least} as many relations as its
child, so $\degop(p)\ge\degop(x)$. The reconstruction is therefore a greedy hunt: attach each node
to the \emph{minimum-degree} neighbour that is still at least as connected as it is.

\solstep{Key Observations}
\begin{observation}
Sorting nodes by degree means when you look at $x$, every possible parent is \emph{later} in the
list. The first later node adjacent to $x$ (smallest such degree) is the closest ancestor -- the
tightest fit, exactly as \secref{c25-degree} reads structure off degrees.
\end{observation}
\begin{observation}
If $x$ and its chosen parent $y$ have the \emph{same} degree, either could be the other's parent
without changing any pair count -- so the tree is not unique: return $2$ if a valid tree exists.
This is the multiplicity signal (\propref{c25-prufdeg} counts a kindred ``swap'' freedom).
\end{observation}
\begin{observation}
Containment is the validity test: for every neighbour $z$ of $x$, $z$ must be adjacent to $y$ (or be
$y$/$x$ itself). A failure means no rooted tree can make $y$ the parent of $x$ -- return $0$. It is
the transitivity the certificate of \thmref{c25-cert} silently assumes.
\end{observation}

\solstep{Think It Through With Me}
\begin{thinkbox}
\deadend{Treat \code{pairs} as tree edges and run a union-find.}{They are \emph{not} edges -- a pair
can join a node to a grandparent. Union-find would happily accept a triangle $\{1,2,3\}$ as a
2-edge tree and miss that two of those are the same ancestor chain. The relation is ancestry, not
adjacency.}
\deadend{Try every root and every parent assignment.}{Exponential. Degrees already order the
search: a parent has degree $\ge$ its child, so a single increasing-degree sweep fixes the candidate
for each node in one pass.}
\aha{Sort by degree; for each node find the first adjacent higher-degree node as its parent, flag an
equal-degree tie as \emph{multiple}, and verify the parent contains all the node's neighbours. Count
roots; more than one root $\to 0$, a tie $\to 2$, else $1$.}
\end{thinkbox}

\solstep{The Approach Ladder}
\begin{center}\renewcommand{\arraystretch}{1.3}
\begin{tabular}{@{}c p{0.235\linewidth} p{0.30\linewidth} l c@{}}
\toprule
\textbf{\#} & \textbf{Approach} & \textbf{Idea} & \textbf{Time} & \textbf{Verdict} \\
\midrule
1 & \lv{3} pairs-as-edges union-find & treat each pair as a tree edge        & $O(m\,\alpha)$ & \textbf{WA} \\
2 & \lv{4} root + parent enumeration & fix root, recurse on ancestor sets    & exponential    & \tle \\
3 & \lv{5} degree-sorted hunt + containment & nearest higher-degree adjacent parent & $O(n^3)$ & \ac \\
\bottomrule
\end{tabular}\end{center}
Rung 1 is the tempting trap -- the pairs are ancestry, not edges, so a cycle-free count means
nothing (Rung 1's failure is precisely why \secref{c25-identify}'s test needs the \emph{right}
edges). Rung 2 is correct but explodes. Rung 3 uses the degree order to pin one parent per node and
the containment check to certify it.

\solstep{Solution}
Build an adjacency \emph{matrix} $\mathit{adj}[\cdot][\cdot]$ and neighbour lists from \code{pairs};
self-loops $\mathit{adj}[v][v]$ are set so the containment test reads uniformly. Collect the nodes
that actually appear, sort them by increasing degree. For each node $x$ (at position $i$), scan
$j>i$ for the first node $y$ adjacent to $x$ -- that is $x$'s parent candidate. If none is found,
$x$ is a root (increment the root count). Otherwise, if $\degop(x)=\degop(y)$ flag \emph{multiple},
and verify every neighbour $z$ of $x$ is adjacent to $y$; a miss returns $0$. At the end, more than
one root returns $0$; otherwise a \emph{multiple} flag returns $2$, else $1$
(\coderef{c25-summit}).

\solstep{Proof of Correctness}
\begin{proofbox}[Nearest higher-degree parent, containment, and the tie count]
\emph{Lemma 1 (a parent has degree at least its child).} If $p$ is $x$'s parent, every node related
to $x$ is related to $p$ (transitivity of ancestry), and $p$ is additionally related to $x$, so
$\degop(p)\ge\degop(x)+1>\degop(x)$ unless $p$ and $x$ have identical relation sets, in which case
$\degop(p)=\degop(x)$. Hence scanning higher-or-equal degree positions only can find $x$'s parent.
\emph{Lemma 2 (the first adjacent one is the closest ancestor).} Among the nodes adjacent to $x$
with degree $\ge\degop(x)$, the minimum-degree one is the deepest possible ancestor -- any other
adjacent node of larger degree is an ancestor of it, not of $x$ directly. Choosing it loses no
valid tree. \emph{Lemma 3 (containment is necessary and sufficient).} $y$ can be $x$'s parent iff
$y$ is adjacent to every neighbour of $x$: \emph{necessary} by transitivity; \emph{sufficient}
because attaching $x$ under such a $y$ preserves exactly the given pairs (no pair is created -- $y$
already relates to all of $x$'s relations -- and none is lost). \emph{Multiplicity and existence.}
If some $x$ shares a degree with its parent $y$, the two are interchangeable (equal $\degop$, equal relation
sets equal), so $>1$ trees exist, flagged by \emph{multiple}; if two nodes have no parent they are
two roots, impossible in one tree ($0$); if a containment fails, no parent works ($0$). Otherwise
the greedy parents are forced, giving exactly one tree. Returning $0/2/1$ by these tests is
therefore correct.
\end{proofbox}

\solstep{Implementation}
\begin{lstlisting}[caption={Number of Ways to Reconstruct a Tree (LeetCode 1719) -- degree-sorted parent hunt + containment}, label={lst:c25-summit}]
#include <bits/stdc++.h>
using namespace std;

class Solution {
public:
    int checkWays(vector<vector<int>>& pairs) {
        int N = 501;                                       // node labels are <= 500
        vector<vector<char>> adj(N, vector<char>(N, 0));   // adjacency as a bit matrix
        vector<vector<int>> nbr(N);
        for (auto& pr : pairs) {
            int A = pr[0], B = pr[1];
            adj[A][B] = adj[B][A] = 1;
            nbr[A].push_back(B); nbr[B].push_back(A);
        }
        vector<int> active;
        for (int id = 1; id < N; ++id)
            if (!nbr[id].empty()) { active.push_back(id); adj[id][id] = 1; }  // self-loop: uniform test
        sort(active.begin(), active.end(),
             [&](int a, int b) { return nbr[a].size() < nbr[b].size(); });    // increasing degree

        bool multiple = false;
        int roots = 0;
        for (int i = 0; i < (int)active.size(); ++i) {
            int x = active[i], j = i + 1;
            while (j < (int)active.size() && !adj[x][active[j]]) ++j;         // first adjacent, >= degree
            if (j < (int)active.size()) {
                int y = active[j];
                if (nbr[x].size() == nbr[y].size()) multiple = true;          // swap -> >1 tree
                for (int z : nbr[x]) if (!adj[y][z]) return 0;                // transitivity fail -> none
            } else {
                ++roots;                                                     // x has no parent -> root
            }
        }
        if (roots > 1) return 0;                                             // two roots impossible
        return multiple ? 2 : 1;
    }
};
\end{lstlisting}

\solstep{Trace}
\begin{tracebox}[Worked example -- \code{pairs = [[1,2],[2,3],[1,3]]}]
All three nodes appear; degrees $\degop(1)=\degop(2)=\degop(3)=2$. After self-loops,
$\mathit{adj}$ is complete on $\{1,2,3\}$. Sorted by degree (tie, order $1,2,3$).
\begin{itemize}[nosep,leftmargin=1.3em]
  \item $x=1$: first adjacent later node is $2$; $\degop(1)=\degop(2)$ so \emph{multiple}$=$true;
        neighbours of $1$ are $\{2,3\}$, both adjacent to $2$ \checkmark.
  \item $x=2$: first adjacent later node is $3$; tie keeps \emph{multiple} true; neighbours
        $\{1,3\}$ both adjacent to $3$ \checkmark.
  \item $x=3$: no later node, so $3$ is the \textbf{root}; $\mathit{roots}=1$.
\end{itemize}
One root and a tie $\Rightarrow$ return $\mathbf{2}$. Indeed the three pairs are the ancestor
relation of a \emph{chain} of length $3$, and either of the two non-root orders ($1\!\to\!2\!\to\!3$
or $3\!\to\!2\!\to\!1$, plus rotations) is valid -- more than one tree. \checkmark
\end{tracebox}
\begin{tracebox}[A quick ``none'' case -- \code{pairs = [[1,2],[2,3]]}]
$\degop=[0,1,2,1]$ for nodes $1,2,3$ (node $1,3$ leaves, node $2$ centre). Sorted: $1,3,2$. $x=1$
attaches to $2$ (adjacent), no tie, neighbours of $1=\{2\}$ contained in $2$ \checkmark. $x=3$
attaches to $2$ \checkmark. $x=2$ is last $\to$ root, $\mathit{roots}=1$. No tie, no containment
break $\Rightarrow$ return $\mathbf{1}$: the unique tree is $2$ as root with children $1,3$.
\end{tracebox}

\solstep{Complexity}
\begin{complexitybox}
Time $\bigO{n^3}$ in the worst case, where $n\le500$ is the number of distinct nodes: the outer
sweep is $n$, the parent scan is $O(n)$, and the containment loop inspects up to $n$ neighbours --
about $500^3=1.25\cdot10^{8}$ elementary operations at the extreme, and $\bigO{m}$ to build the
matrix ($m\le2\cdot10^{5}$). Space $\bigO{n^2}$ for the adjacency matrix ($501^2$ bytes). The
$n\le500$ label bound is what makes a dense $\mathit{adj}[\cdot][\cdot]$ affordable here -- contrast
\secref{c25-parent}, where an edge list would blow up at $n$ large.
\end{complexitybox}

\solstep{Verification and Edge Cases}
\begin{center}\renewcommand{\arraystretch}{1.2}
\begin{tabular}{@{}p{0.30\linewidth} p{0.62\linewidth}@{}}
\toprule
\textbf{Case} & \textbf{What the code does} \\
\midrule
Star (one node in every pair) & that node has max degree, becomes the single root; leaves tie-free $\to 1$ \\
Perfect chain & every consecutive degree differs by one; no tie, one root $\to 1$ \\
Two nodes same relation set & equal-degree tie fires $\to$ \emph{multiple} $\to 2$ \\
Contradictory pairs & a containment check finds a neighbour $z$ with $\mathit{adj}[y][z]=0$ $\to 0$ \\
Two components (two roots) & \code{roots}$>1$ at the end $\to 0$ \\
\bottomrule
\end{tabular}\end{center}

\solstep{Pitfalls and Variants}
\begin{warningbox}
Three classic bugs: (i) reading \code{pairs} as tree \emph{edges} and running \corref{c25-istree}'s
$n-1$-count test -- they are \emph{ancestry} pairs, so the count is meaningless (Rung 1's WA).
(ii) Forgetting the equal-degree tie: a valid tree with a swappable parent must report $2$, not $1$.
(iii) Skipping the \code{roots}$>1$ guard -- a forest of two ``trees'' has two parentless nodes and
must be $0$. Variant -- if the problem asked to \emph{build} one tree rather than count, the same
greedy parents give it directly; and the ``distinct labels $\le500$'' bound is what lets a dense
matrix stand in for the adjacency lists of \secref{c25-parent}.
\end{warningbox}

\begin{insightbox}[Key takeaways -- three things to remember]
\begin{enumerate}[nosep, leftmargin=1.6em]
  \item \emph{Partial data still pins structure.} Ancestor pairs, though not edges, order a unique
        parent per node once you sort by degree -- the degree-reading skill of \secref{c25-degree}.
  \item \emph{Transitivity is the certificate here.} Containment ($y$ relates to all of $x$'s
        relations) plays the role the union-find rejection plays in \secref{c25-identify}.
  \item \emph{Multiplicity $\iff$ an interchangeable slot.} Equal-degree adjacent nodes can swap,
        exactly as \propref{c25-prufdeg}'s count is $>1$ whenever two degrees repeat.
\end{enumerate}
\end{insightbox}

\begin{hardeningbox}
\begin{itemize}[nosep, leftmargin=1.4em]
  \item \textbf{Larger label range}: replace the $\bigO{n^2}$ \code{char} matrix with a hash set of
        neighbours and test containment by set-subtraction -- same logic, $\bigO{m}$ memory, so the
        $n\le500$ ceiling lifts.
  \item \textbf{Return a real tree}: keep the chosen $y$ as \code{parent[x]} during the sweep; the
        greedy is already a reconstruction, closing the loop with \secref{c25-parent}.
  \item \textbf{This is counting-in-disguise}: the ``how many'' $\{0,1,2\}$ answer mirrors
        \propref{c25-prufdeg}'s exact count $(n-2)!/\prod(\degop-1)!$; the full Pr\"ufer bijection
        that turns such counts into closed forms is \chref{c32}.
\end{itemize}
\end{hardeningbox}

\chaptersection{Tool Kit Inventory}

The machine assembled so far -- cumulative through Chapter~6. Chapters~1--5's rows are carried
forward; Chapter~6's rows are appended. Pointers only, never prose.

\begin{center}
\small
\renewcommand{\arraystretch}{1.25}
\begin{tabular}{@{}p{0.24\linewidth} l c p{0.27\linewidth} p{0.16\linewidth}@{}}
\toprule
\textbf{Tool} & \textbf{Since} & \textbf{Level} & \textbf{Trigger phrase} & \textbf{Where} \\
\midrule
\multicolumn{5}{l}{\itshape Chapters 1--5 -- carried forward} \\
Adjacency list + parent-guard DFS & 1.1 & \lv1 & edge list given, not pointers & \coderef{c01-build} \\
Three-slot DFS (pre/in/post) & 3.1 & \lv1 & a node sequence in a named order & \coderef{c03-dfs} \\
Root: read parent, depth, size, height & 4.2 & \lv3 & undirected edges, need above/below & \coderef{c04-root} \\
Stamp $\tin/\tout$; $O(1)$ ancestor test & 5.1 & \lv2 & ``is $u$ above $v$?'' / DFS-order listing & \coderef{c24-stamp} \\
Flatten: subtree $=$ range $[\tin,\tin{+}\size)$ & 5.2 & \lv3 & many whole-subtree queries & \coderef{c24-flatten-build} \\
\midrule
\multicolumn{5}{l}{\itshape Chapter 6 -- new this chapter} \\
Rebuild adj, read root, certify & 6.1 & \lv3 & edge list / parent array, valid rooted tree? & \coderef{c25-build} \\
Degree-sum gate + leaf-peel realisation & 6.2 & \lv3 & ``build/count a tree from degrees'' & \coderef{c25-degbld} \\
Union-find pre-look (tree? components?) & 6.3 & \lv2 & ``is it a tree / how many provinces'' & \coderef{c25-uf} \\
Optimal guess on a 3-way oracle & 6.4 & \lv2 & ``find a hidden value, fewest questions'' & \coderef{c25-guess} \\
Degree-greedy parent hunt + containment & Summit & \lv5 & ``reconstruct from ancestor pairs'' & \coderef{c25-summit} \\
\bottomrule
\end{tabular}
\end{center}

\chaptersection{Related Topics}

\begin{itemize}[nosep, leftmargin=1.4em]
  \item \chref{c03}, \chref{c04} -- backward: rebuilding from a \emph{pair of traversals}
        (in+pre) lives in \chref{c03}; this chapter rebuilds from \emph{arbitrary} data, and the
        mechanical adjacency build / parent-guarded rooting is \secref{c04-rooting} -- we add only
        the validity and root-detection decisions.
  \item \chref{c24} -- forward/backward: 6.4 \emph{uses} the distance/ancestor oracle ($\tin/\tout$,
        $\dist$, LCA) as a black box to answer queries about a hidden tree; it never re-derives it.
  \item \chref{c28} -- forward: the full Disjoint-Set-Union structure, union by rank and path
        compression live there (its greedy use in MST is Ch.~31, \chref{c17}); \secref{c25-identify} is only the one-pass pre-look.
  \item \chref{c32} -- forward: the Pr\"ufer-code bijection and Cayley's $n^{\,n-2}$ proof; here the
        code is a \emph{teaser} that explains why a degree sequence pins a counting problem
        (\propref{c25-prufdeg}).
  \item \chref{c11}, \chref{c21} -- LCA and Mo-on-trees both \emph{assume} a rebuilt, rooted tree
        -- the object this chapter produces and certifies.
\end{itemize}

\begin{lastword}[Last Word -- Chapter 6]
\emph{``Every earlier chapter walked a tree you were \textbf{given}; this chapter asked whether a
tree can be summoned back from the \textbf{marks} it left -- and proved the summoning was honest.''}

\medskip
Chapters~3--5 taught you to traverse, store, and flatten a tree that was already drawn for you.
This chapter inverted the arrow. You were handed a pile of pairs and asked to \emph{rebuild} the
shape (\secref{c25-parent}), handed a list of degrees and asked whether \emph{any} tree wears those
numbers (\secref{c25-degree}), handed an edge stream and asked \emph{is this even a tree}
(\secref{c25-identify}), and finally handed nothing but an oracle and asked to \emph{elicit} the
answer one question at a time (\secref{c25-query}). Four shadows -- edge list, parent array, degree
sequence, distance answers -- all of the \emph{same} $T$ in \figref{c25-running}, and each enough to
rebuild it.

\medskip
Carry three reflexes forward:
\begin{enumerate}[nosep, leftmargin=1.6em]
  \item \emph{Certify with two legs, never one.} $n-1$ edges \emph{and} connected, or $n-1$ edges
        \emph{and} acyclic -- any two of the certificate force the third (\thmref{c25-cert}); the
        count alone is a trap a forest happily passes.
  \item \emph{Degrees are data.} $\sum\degop=2(n-1)$ decides realisability, $(n-2)!/\prod(\degop-1)!$
        counts the trees, and the Summit's parent hunt runs entirely on degree order
        (\thmref{c25-degreal}, \propref{c25-prufdeg}).
  \item \emph{A good question halves the world.} Against an oracle, choose the query whose answers
        leave equally small sub-cases; $\lceil\log_2\rceil$ is not a convenience, it is a
        lower bound (\thmref{c25-logn}).
\end{enumerate}

\medskip
And three habits, drilled until boring: write $T$'s parent array, degree sequence, and Pr\"ufer
code by hand and check $\sum\degop=2(n-1)$ before you trust a rebuild; keep straight the
\emph{undirected} cycle test (union-find) from the \emph{directed} one (topological sort) -- they
answer different questions and each fails the other; and never answer an oracle by a linear scan
when a midpoint halves the candidate set.

\medskip
\textbf{Next:} \chref{c05} keeps building on the rebuild instinct but installs an \emph{order} inside
the tree -- a Binary Search Tree is a tree that knows its own inorder, so search, insert, delete and
rank queries all collapse into a single root-to-leaf descent, and the degeneration of that descent is
the cliffhanger that opens the balancing chapter.

\begin{center}\textcolor{softpurple}{\rule{0.4\linewidth}{0.6pt}\quad $\blacksquare$\quad \rule{0.4\linewidth}{0.6pt}}\end{center}
\end{lastword}


\closechapter
